\documentclass[11pt]{amsart}
\usepackage{amssymb,amsthm,mathtools}
\usepackage{array,booktabs,longtable}
\usepackage{url}
\usepackage[colorlinks=true,linkcolor=blue,citecolor=blue,urlcolor=blue]{hyperref}

\allowdisplaybreaks[1]
\raggedbottom
\setlength{\emergencystretch}{2em}
\hyphenation{Aron-szajn Orn-stein Uhlen-beck}

% the macros of paper.tex
\newcommand{\R}{\mathbb{R}}
\newcommand{\Sph}{\mathbb{S}}
\newcommand{\Tr}{\operatorname{Tr}}
\newcommand{\dd}{\,\mathrm{d}}
\newcommand{\FC}[1]{\mathcal{F}_{#1}}
\newcommand{\Ct}{\widetilde{C}_t}
\newcommand{\taut}{\widetilde{\tau}_t}
\newcommand{\ip}[2]{\langle #1,#2\rangle}
\newcommand{\hatf}[1]{\widehat{#1}}
\newcommand{\tp}{\top}

% structure of the commentary
\newcommand{\src}[1]{\par\bigskip\noindent\textbf{Lines #1.}\enspace}
\newcommand{\content}{\par\smallskip\noindent\textit{Content.}\enspace}
\newcommand{\why}{\par\smallskip\noindent\textit{Why it holds.}\enspace}
\newcommand{\verified}{\par\smallskip\noindent\textit{Verified.}\enspace}
\newcommand{\notes}{\par\smallskip\noindent\textit{Notes.}\enspace}
\newcommand{\tA}{\textsf{A}}
\newcommand{\tS}{\textsf{S}}
\newcommand{\tN}{\textsf{N}}
\newcommand{\tL}{\textsf{L}}
\newcommand{\tB}{\textsf{B}}
\newcommand{\file}[1]{\texttt{#1}}
\newcommand{\PS}[1][]{[PSWPN\ifx\relax#1\relax\else, #1\fi]}
\makeatletter\g@addto@macro\UrlBreaks{\do\-}\makeatother
\newcolumntype{L}[1]{>{\raggedright\arraybackslash}p{#1}}

\begin{document}

\title[Commentary on the manuscript]{Line-by-line commentary and submission audit for the manuscript
``Koopman invariance of Gaussian RKHS$\text{s}$ for linear stochastic differential equations''}
\date{24 September 2026}

\begin{abstract}
This document goes through \file{paper.tex} (1021 source lines, 19 pages when compiled) in source order. For each
block of source lines it states what the lines say and derives every mathematical claim with the intermediate steps
written out. It also records how each claim was checked and notes caveats. Section~\ref{sec:summary} summarizes the
audit, including the corrections made during it. Section~\ref{sec:ready} answers whether the manuscript is ready for
submission and to which journal it should go.
\end{abstract}

\maketitle
\tableofcontents

%======================================================================
\section{How to read this document}\label{sec:howto}
%======================================================================

Line numbers refer to \file{paper.tex} as committed together with this commentary. Each block of lines is discussed
under up to four headings:
\begin{itemize}
\item \emph{Content}: what the lines say;
\item \emph{Why it holds}: a derivation, with all intermediate steps;
\item \emph{Verified}: how the claim was checked;
\item \emph{Notes}: caveats, and corrections made during this audit.
\end{itemize}
The verification tags are the following.
\begin{itemize}
\item[\tA] Re-derived by hand. The derivation is written out in this document.
\item[\tS] Confirmed in exact arithmetic by \file{verify\_symbolic.py} (SymPy, 52 checks in sections S1--S5 of its
output).
\item[\tN] Confirmed in floating-point arithmetic by \file{verify.py} (38 checks) or \file{make\_figure.py}. The item
of Appendix~A of the paper is named.
\item[\tL] Compared with the cited source; the location is named. For \PS\ the source is the arXiv PDF supplied with
the task, \emph{arXiv:2405.14429v1, 23 May 2024}. Quotations from it are verbatim, except that its mathematical
notation has been re-typeset.
\item[\tB] Build checks: \file{pdflatex} (TeX Live 2023) and a KaTeX parse of all 960 mathematical snippets of
\file{paper.tex}.
\end{itemize}
The notation is that of the paper. We write $P\le Q$ if $Q-P$ is positive semidefinite (PSD), $|x|$ for the Euclidean
norm and $\|M\|$ for the spectral norm. The Fourier transform is $\hatf g(\xi)=\int g(x)e^{-i\xi\cdot x}\dd x$, and
\begin{gather*}
w_C(\xi)=e^{\frac14\xi^\tp C^2\xi},\qquad\kappa_C=\pi^{d/2}\det C,\\
D_t=C^2+4\Sigma(t)-e^{At}C^2e^{A^\tp t},\qquad\Lambda_C=4BB^\tp-(AC^2+C^2A^\tp).
\end{gather*}
Further, $C_t^2=e^{-At}(C^2+2\Sigma(t))e^{-A^\tp t}$ and $\tau_t=\det C/\det(C^2+2\Sigma(t))^{1/2}$ are the objects
of \PS, while $\Ct^{\,2}=e^{-At}(C^2+4\Sigma(t))e^{-A^\tp t}$ and $\taut=\det C/\det(C^2+4\Sigma(t))^{1/2}$ are those
of the paper. For Hurwitz $A$ we put $P_C=C^2+4\Sigma$. Equation, theorem and section numbers such as ``(3.4)'' or
``Theorem 4.2'' refer to the paper; the sections of this commentary are cited as ``Section~\ref{sec:summary}'' and so
on.

%======================================================================
\section{Summary of the audit}\label{sec:summary}
%======================================================================

\subsection{Result}
I found no mathematical error in any statement or proof of the manuscript. Every displayed formula and every proof step
was re-derived; the derivations are in Sections~\ref{sec:front}--\ref{sec:app}. The following checks were run on 24
September 2026.

\medskip
{\small
\begin{longtable}{L{140pt}L{196pt}}
\toprule
Check & Result\\
\midrule
\endhead
\file{verify.py} (Appendix A, items 1--6 and 8--10) & 38 of 38 checks pass; the output is byte-identical to the
committed \file{verify\_output.txt}.\\
\file{make\_figure.py} (Figure 1, item 7) & The interval property and the tail bound hold for all 320 values of
$\beta$. The printed results are identical to \file{figure\_output.txt}, and the regenerated figure differs from the
committed one only in its creation date.\\
\file{verify\_symbolic.py} (new) & 52 of 52 exact checks pass.\\
KaTeX parse of \file{paper.tex} & 960 mathematical snippets, 0 errors.\\
\file{pdflatex}, two passes & 0 errors, 0 warnings, 0 overfull or underfull boxes, 19 pages.\\
Citations of \PS & Every numbered reference matches arXiv v1: Theorem 2.1 and its proof, (2.3), Remark 2.2(c),
Corollary 2.3, Lemma 3.1, Proposition 3.3, (4.1), Proposition 4.1 and its proof, Remark 4.2.\\
\bottomrule
\end{longtable}
}

\subsection{What is proved, what is computed}
All theorems, lemmas, propositions and corollaries are proved analytically. The same holds for every statement in
Examples 6.1 and 6.2, with two exceptions: the curve $t_1(\beta)$, and the statement that $\{t\in(0,40]:D_t\ge0\}$ is
an interval for each of the 320 sampled values of $\beta$. Both are numerical and are labelled as such in the paper.
The error magnitudes and the Monte Carlo value in Appendix A are numerical. No proof depends on a numerical result.

\subsection{Corrections made during the audit}\label{sec:corrections}
The line numbers refer to the revised \file{paper.tex}.
\begin{itemize}
\item[C1] \emph{Notation} (lines 51, 89). The Euclidean norm of a vector was written $\|\cdot\|$, although the
Notation paragraph (lines 182--186) reserves $\|\cdot\|$ for the spectral norm and $|\cdot|$ for vectors. The abstract
now writes the quadratic form $(x-y)^\tp C^{-2}(x-y)$, and line 89 uses $|\cdot|$.
\item[C2] \emph{Attribution} (lines 150--153, 758--764, 788--789, 889--899). The formula
$K^tk^C_z=\tau_t\,k^{C_t}_{e^{-At}z}$ is established in the \emph{proof} of \PS[Proposition 4.1]; the statement only
bounds the norm of $K^t\colon H_C\to H_{C_t}$ by $\tau_t^{1/2}$. All citations now say so, and Remark 7.2 derives the
formula from (3.2) and (2.3). The paper is therefore self-contained on this point.
\item[C3] \emph{Self-containedness} (lines 518--520). The proof of Corollary 4.3 quoted the derivative $\dot f_x(0)$
from \PS. It now derives it in one line.
\item[C4] \emph{An unproved claim} (lines 157--160 and 765--769). Remark 7.2 said that non-invariance is witnessed
``only'' by functions whose Fourier transforms decay slowly, and this was not proved. It is replaced by the statement
that is proved: a finite kernel expansion $f$ satisfies $|\hatf f|^2w_C\le c\,w_C^{-1}$, so it is never a witness
once $C^2+D_t>0$. The introduction was adjusted accordingly.
\item[C5] \emph{A missing hypothesis} (lines 636--637). Remark 5.2 spoke of ``the ellipsoid of $\Sigma$'', but
Theorem 5.1 does not assume controllability, so $\Sigma$ may be singular. The remark now says ``whenever $\Sigma>0$''.
\item[C6] \emph{Imprecise explanation} (lines 851--855). Example 7.6 attributed condition (1.5) to ``Gaussian
kernels''. It now names the mechanism: the Gaussian decay of the spectral density competes with the factor
$\exp(-\eta^\tp\Sigma(t)\eta)$. It also cites the related Sobolev results \cite{Koehne,Hertel}.
\item[C7] \emph{Ambiguity} (lines 58--60). In the abstract, ``the criterion says'' could refer to the fixed-time
criterion. The sentence now refers explicitly to invariance for all $t\ge0$.
\item[C8] \emph{Precision} (lines 127--128 and 344--345). The phrase ``for the bounded continuous functions in $H_C$''
suggested that only some elements of $H_C$ are bounded and continuous; in fact all are. Also, ``bounded and
continuous, since $|f(x)|\le\|f\|_C$'' justified only boundedness; continuity is now attributed to Lemma 2.1(b).
\item[C9] \emph{Precision} (lines 142--143). The Fourier membership criterion in ``Idea of the proof'' now states its
hypotheses: $f$ continuous, real-valued and in $L^2$.
\item[C10] \emph{A wrong norm} (lines 873--885 and 889--892). Appendix A, item 2 wrote $\|t^{-1}D_t-\Lambda_C\|$, which
by the Notation paragraph is the spectral norm, whereas \file{verify.py} measures the largest entry. The appendix now
defines $|M|_{\max}=\max_{i,j}|M_{ij}|$ and states items 1, 2 and 4 in this norm.
\item[C11] \emph{Completeness} (lines 871--873 and 920--925). Appendix A now describes the third, exact-arithmetic
script.
\item[C12] \emph{Literature} (lines 77, 797--798, 854--855 and the bibliography). The audit added:
\begin{itemize}
\item Hertel, Philipp, Schaller and Worthmann \cite{Hertel}, a follow-up of \PS\ on stochastic kernel EDMD;
\item Ikeda, Ishikawa and Sawano \cite{IIS}, on composition operators on RKHSs of analytic kernels.
\end{itemize}
Köhne et al.\ was updated from preprint to \emph{SIAM J.\ Appl.\ Dyn.\ Syst.} 24 (2025). Mauroy et al.\ gained its
subtitle and series volume. A note records that numbered statements of \PS\ follow arXiv v1.
\item[C13] \emph{Dead code} (line 10). The hyphenation exception \texttt{Mat-ern} was removed. It can never match the
word ``Mat\'ern'', because the accent command interrupts the word.
\item[C14] \emph{Submission items} (lines 42--43, 927--935, 938--939). The audit added a commented AI-use statement
and a reminder next to the commented author block. The data availability sentence now reads ``No external data were
used'', which is accurate; before, it said that no data were generated.
\item[C15] \emph{Style} (lines 70--71, 91). The first sentence of the introduction was split for grammar, and ``(real)
RKHS'' is now explicit.
\end{itemize}

\subsection{Novelty check}\label{sec:novelty}
Web searches on 24 September 2026 found no publication that answers the question of \PS[Remark 2.2(c), Remark 4.2]
\cite{src-survey,src-hertel}. According to the search engine's summary, the July 2026 survey arXiv:2607.01819 still
describes the question as open. I could not open the PDF, because the network policy of this environment blocks
\url{arxiv.org}; the conclusion rests on search results, not on reading the survey. The closest works are:
\begin{itemize}
\item \cite{Hertel}, on Sobolev spaces for stochastic dynamics, via conditions on a Radon--Nikodym density;
\item \cite{Koehne}, on Wendland (Sobolev) native spaces for deterministic dynamics;
\item \cite{IIS}, on composition operators.
\end{itemize}
None of them treats Gaussian RKHSs of linear SDEs. A search cannot prove a negative. Before submission the author
should repeat the search on arXiv and MathSciNet/zbMATH, which were not reachable from here.

%======================================================================
\section{Front matter (lines 1--64)}\label{sec:front}
%======================================================================

\src{1--10}
\content Document class \texttt{amsart} at 11\,pt. The packages are \texttt{amssymb} (blackboard bold),
\texttt{amsthm} (theorem environments), \texttt{mathtools} (a superset of \texttt{amsmath}: \texttt{align*},
\texttt{gather*}, \texttt{smallmatrix}), \texttt{graphicx} (Figure 1) and \texttt{hyperref} (links). The settings
work as follows.
\begin{itemize}
\item \verb|\numberwithin| numbers equations as (section.number).
\item \verb|\allowdisplaybreaks[1]| permits page breaks inside multi-line displays.
\item \verb|\raggedbottom| lets pages end at their natural height.
\item \verb|\emergencystretch| adds a third line-breaking pass with extra stretch.
\item \verb|\hyphenation| lists exceptions for proper names.
\end{itemize}
\verified \tB: the log has no overfull or underfull boxes; the last three settings were introduced to achieve this.
\notes Journals reformat with their own class (Section~\ref{sec:ready}).

\src{12--19}
\content Theorem-like environments that share one counter numbered within sections. So Theorem 4.1, Theorem 4.2,
Corollary 4.3, Theorem 4.4, Remark 4.5 and Corollary 4.6 are consecutive. Examples use the definition style (upright
text), and remarks use the remark style.

\src{21--30}
\content The macros are the following.
\begin{itemize}
\item \verb|\R| is $\R$, \verb|\Sph| is $\Sph$ (the sphere $\Sph^{d-1}$), \verb|\Tr| is the trace.
\item \verb|\dd| is the differential with a thin space.
\item \verb|\FC{C}| is the space $\FC{C}$ of Fourier transforms.
\item \verb|\Ct| is $\Ct$ and \verb|\taut| is $\taut$.
\item \verb|\ip{f}{g}| is $\ip fg$, \verb|\hatf| is the wide hat, and \verb|\tp| is the transpose sign.
\end{itemize}
\verified \tB: KaTeX parses all 960 mathematical snippets with these macro definitions.

\src{32--37}
\content Title and short title. amsart sets titles and running heads in capitals with \verb|\uppercasenonmath|, which
does not change mathematics. Writing the plural as \verb|RKHS$\text{s}$| therefore keeps it as ``RKHSs'' instead of
``RKHSS''.

\src{38--43}
\content The author block is commented out and must be completed by the submitting author. Lines 42--43 remind the
author to activate the AI-use statement (lines 927--935).
\notes This is a blocking item for submission (Section~\ref{sec:ready}).

\src{45--47}
\content MSC 2020 classification and keywords. Primary codes:
\begin{itemize}
\item 46E22, Hilbert spaces with reproducing kernels;
\item 47D07, Markov semigroups and applications to diffusion processes.
\end{itemize}
Secondary codes:
\begin{itemize}
\item 60H10, stochastic ordinary differential equations;
\item 47B33, linear composition operators;
\item 42B10, Fourier and Fourier--Stieltjes transforms and other transforms of Fourier type.
\end{itemize}
\notes The code 47B32 (linear operators in reproducing-kernel Hilbert spaces) would also fit. Adding it is optional.

\src{49--62}
\content The abstract. Sentence by sentence:
\begin{enumerate}
\item The result of \PS: sufficiency of $\frac12(AC^2+C^2A^\tp)\le BB^\tp$ under ``$A$ Hurwitz, $(A,B)$
controllable''. \tL: \PS[Theorem 2.1] and the standing assumptions of \PS[Section 2].
\item ``They asked whether this condition is also necessary.'' \tL: \PS[Remark 2.2(c)] reads ``It is not clear whether
the condition (2.3) is also necessary for Koopman invariance of $H_C$.'' \PS[Remark 4.2] reads ``It is not clear
whether the statement of Theorem 2.1 is actually an equivalence.''
\item The answer and the sharp condition: Theorem 1.1(b), which is Theorem 4.4, and Corollary 1.2.
\item ``$K^t$ maps $H_C$ isomorphically onto the Gaussian RKHS with matrix $\Ct$'': Theorem 4.2. The map is
$\taut^{1/2}$ times a unitary operator.
\item The fixed-time criterion and the norm: Corollary 4.3 and Theorem 4.1.
\item The Hurwitz case: Theorem 5.1(b).
\item The method: Lemmas 2.1 and 3.2.
\end{enumerate}
\notes Corrections C1 and C7 apply here.

%======================================================================
\section{Paper Section 1: Introduction (lines 67--186)}\label{sec:intro}
%======================================================================

\src{70--79}
\content The Koopman operator of a stochastic system is $(K^tf)(x)=\mathbb E[f(X_t)\,|\,X_0=x]$. It is linear, it can
be learned from data (EDMD and its kernel variants), and error bounds for kernel methods profit from the invariance of
the RKHS.
\why Linearity: $K^t(\alpha f+\beta g)=\alpha K^tf+\beta K^tg$ by linearity of conditional expectation, whatever the
dynamics.
\verified \tL: the bibliographic data are checked in Section~\ref{sec:bib}. The remaining statements are descriptive.

\src{81--103}
\content The setting of \PS.
\begin{itemize}
\item The SDE (1.1).
\item The kernel $k^C$ and its real RKHS $H_C$.
\item The conditional covariance $\Sigma(t)$ in (1.2).
\item The standing assumptions: $m\le d$, $(A,B)$ controllable, $A$ Hurwitz.
\item The sufficient condition (1.3), under which $\|K^tf\|_C\le e^{-t\Tr A/2}\|f\|_C$.
\end{itemize}
The open question is whether (1.3) is also necessary.
\why That $\Sigma(t)$ is the covariance of $X_t$ given $X_0$ is shown in Section~\ref{sec:koop} (lines 336--343).
\verified \tL. \PS[Section 2] assumes ``$m\le d$'', ``$(A,B)$ is controllable and $A$ is Hurwitz''.
\PS[Theorem 2.1] reads: ``Let $C\in\R^{d\times d}$ be a symmetric positive definite matrix such that
$\frac12(AC^2+C^2A^\top)\le BB^\top$. (2.3) Then the RKHS $H_C$ is invariant under each Koopman operator $K^t$,
$t\ge0$, associated with the SDE (2.1), and we have $\|K^tf\|_C\le e^{t\cdot\Tr(-A)/2}\cdot\|f\|_C$.'' Both remarks
cited at line 103 are quoted in Section~\ref{sec:front}.
\notes Corrections C1 and C15 apply here.

\src{105--125}
\content Theorem 1.1. For arbitrary $A$ and $B$ and symmetric positive definite $C$:
\begin{itemize}
\item[(a)] $K^tH_C\subset H_C$ if and only if $e^{At}C^2e^{A^\tp t}\le C^2+4\Sigma(t)$, and then
$\|K^t\|=e^{-t\Tr A/2}$;
\item[(b)] invariance for all $t\ge0$ holds if and only if $\frac12(AC^2+C^2A^\tp)\le2BB^\tp$. If this fails,
invariance fails on a whole interval $(0,\varepsilon)$.
\end{itemize}
\why (a) is Theorem 4.1(a),(b), because $D_t\ge0$ is (1.4). (b) is Theorem 4.4, and the ``interval'' statement is
proved at lines 558--563 (Section~\ref{sec:proofs}).
\notes
\begin{itemize}
\item No standing assumption is needed: $A$ may be unstable, $B$ may vanish, and $m$ is arbitrary.
\item Part (a) shows that the bound of \PS[Theorem 2.1] is \emph{sharp}. Whenever $H_C$ is invariant,
$\sup_{f\ne0}\|K^tf\|_C/\|f\|_C=e^{-t\Tr A/2}$. For Hurwitz $A$ this number exceeds $1$ and grows exponentially in $t$.
\end{itemize}

\src{127--129}
\content $K^tf$ is defined pointwise, which is meaningful because every $f\in H_C$ is bounded and continuous. Under the
assumptions of \PS\ this agrees with the $L^2(\mu)$ definition.
\why See lines 344--345 and Remark 3.1 (Section~\ref{sec:koop}).
\notes Correction C8 applies here.

\src{131--139}
\content Corollary 1.2. Condition (1.3) implies (1.5) but not conversely. The matrices
$A=\begin{pmatrix}-1&5\\0&-1\end{pmatrix}$ and $B=C=I$ separate the two conditions. Moreover, (1.5) is (1.3) with
$C/\sqrt2$ in place of $C$.
\why See the proof at lines 565--572 (Section~\ref{sec:proofs}). The last statement holds because
\[
\tfrac12\bigl(A(C/\sqrt2)^2+(C/\sqrt2)^2A^\tp\bigr)=\tfrac14(AC^2+C^2A^\tp),
\]
so ``$\le BB^\tp$'' for $C/\sqrt2$ is the same as ``$\frac12(AC^2+C^2A^\tp)\le2BB^\tp$'' for $C$.
\verified \tS: S5 (the identity) and S3 (eigenvalues $\frac32$ and $-\frac72$ for $\beta=5$). \tN: item 6 ($\beta=5$).
\notes The following observation is not in the paper. Since
\[
k^{C/\sqrt2}(x,y)=\exp\bigl(-2(x-y)^\tp C^{-2}(x-y)\bigr)=k^C(x,y)^2,
\]
the sharp condition for $H_C$ is exactly the sufficient condition of \PS\ for the \emph{squared} kernel $(k^C)^2$. It
could be added as a one-line remark.

\src{141--146}
\content Idea of the proof: a Fourier description of $H_C$; the Fourier form of $K^t$; comparison of two Gaussian
weights; differentiation at $t=0$.
\why Lemma 2.1, Lemma 3.2, Theorem 4.1 and Theorem 4.4.
\notes Correction C9 applies here.

\src{148--160}
\content ``Where the factor two comes from.'' \PS\ track kernel sections. Averaging a Gaussian bump over the Gaussian
transition law adds $\Sigma(t)$ to its covariance $\frac12C^2$, and this gives $K^tk^C_z=\tau_tk^{C_t}_{e^{-At}z}$.
Membership in $H_C$ is decided by $|\hatf f|^2$, where the smoothing factor appears squared. This turns $2\Sigma(t)$
into $4\Sigma(t)$.
\why \emph{Kernel sections.} By (2.3), $\hatf{k^C_z}(\eta)=\kappa_Ce^{-i\eta\cdot z}w_C(\eta)^{-1}$. By (3.2) with
$\eta=e^{-A^\tp t}\xi$, and since $(e^{-A^\tp t}\xi)\cdot z=\xi\cdot e^{-At}z$,
\begin{align*}
\hatf{K^tk^C_z}(\xi)
&=e^{-t\Tr A}\kappa_Ce^{-i\xi\cdot e^{-At}z}
\exp\bigl(-\tfrac14\xi^\tp e^{-At}C^2e^{-A^\tp t}\xi\bigr)\\
&\qquad\times\exp\bigl(-\tfrac12\xi^\tp e^{-At}\Sigma(t)e^{-A^\tp t}\xi\bigr)\\
&=e^{-t\Tr A}\kappa_Ce^{-i\xi\cdot e^{-At}z}\,w_{C_t}(\xi)^{-1}.
\end{align*}
Also $\hatf{k^{C_t}_{e^{-At}z}}(\xi)=\kappa_{C_t}e^{-i\xi\cdot e^{-At}z}w_{C_t}(\xi)^{-1}$. The two transforms therefore
differ by the factor
\[
\frac{e^{-t\Tr A}\kappa_C}{\kappa_{C_t}}=\frac{e^{-t\Tr A}\det C}{\det C_t}
=\frac{e^{-t\Tr A}\det C}{e^{-t\Tr A}\det(C^2+2\Sigma(t))^{1/2}}=\tau_t ,
\]
where $\det C_t=(\det e^{-At})\det(C^2+2\Sigma(t))^{1/2}=e^{-t\Tr A}\det(C^2+2\Sigma(t))^{1/2}$. Both sides are
continuous $L^2$ functions with the same Fourier transform, hence $K^tk^C_z=\tau_tk^{C_t}_{e^{-At}z}$.

\emph{Squared factor.} For general $f\in H_C$, (3.2) gives
\[
|\hatf{K^tf}(\xi)|^2=e^{-2t\Tr A}|\hatf f(e^{-A^\tp t}\xi)|^2\exp\bigl(-\xi^\tp e^{-At}\Sigma(t)e^{-A^\tp t}\xi\bigr).
\]
The smoothing factor now carries $\Sigma(t)$, not $\frac12\Sigma(t)$. Let $C'$ be a candidate matrix and substitute
$\xi=e^{A^\tp t}\eta$. Then the weight $\exp\bigl(-\eta^\tp\Sigma(t)\eta+\frac14\eta^\tp e^{At}C'^2e^{A^\tp t}\eta\bigr)$
must be compared with $w_C(\eta)$. The two are equal exactly when $e^{At}C'^2e^{A^\tp t}=C^2+4\Sigma(t)$, i.e.\ for
$C'=\Ct$ (Theorem 4.2).

\emph{Loss in \PS.} Since $\Ct^{\,2}-C_t^2=2e^{-At}\Sigma(t)e^{-A^\tp t}\ge0$, the range $H_{\Ct}$ is a subspace of
$H_{C_t}$, and a proper one when $\Sigma(t)\ne0$ (Corollary 4.3). Checking $H_{C_t}\subset H_C$ asks for more than
invariance needs.
\verified \tL. The formula appears in the proof of \PS[Proposition 4.1], whose last display ends with ``hence,
$K^tk^C_z=\tau_t\cdot k^{C_t}_{e^{-At}z}\in H_{C_t}$''. The statement of \PS[Proposition 4.1] says that $K^t$ ``maps
$H_C$ boundedly into $H_{C_t}$ with norm not exceeding $\tau_t^{1/2}$''. The equivalence ``(1.3) $\iff$
$H_{C_t}\subset H_C$ for all $t$'' is in the proof of \PS[Theorem 2.1]: ``We prove that $C_t^2\ge C^2$ for all $t\ge0$
if and only if (2.3) is satisfied'', combined with \PS[Proposition 3.3]. \tS: S2 checks the kernel-section formula by
an $x$-space Gaussian convolution for $d=1$. \tN: item 5 is a Monte Carlo check.
\notes Corrections C2 and C4 apply here.

\src{161--167}
\content For Hurwitz $A$, $\Sigma=\lim_{t\to\infty}\Sigma(t)$ exists and solves $A\Sigma+\Sigma A^\tp=-BB^\tp$. The
two conditions become Lyapunov inequalities for $C^2+4\Sigma$ and for $C^2+2\Sigma$. Invariance is equivalent to
forward invariance of an ellipsoid, and it holds for all large $t$.
\why Existence of the limit: if $A$ is Hurwitz, there are $M\ge1$ and $\alpha>0$ with $\|e^{As}\|\le Me^{-\alpha s}$.
Hence $\|e^{As}BB^\tp e^{A^\tp s}\|\le M^2\|B\|^2e^{-2\alpha s}$ is integrable on $[0,\infty)$. The Lyapunov equation
and the equivalences are Theorem 5.1 (Section~\ref{sec:geom}).

\src{169--173}
\content The Fourier argument extends to translation-invariant kernels (Proposition 7.5). Sobolev (Mat\'ern) spaces
are always invariant (Example 7.6).

\src{175--186}
\content Organization and notation.
\verified \tA: the section list matches the section titles at lines 188, 333, 412, 588, 645, 735 and 868.

%======================================================================
\section{Paper Section 2: Gaussian RKHSs in Fourier variables (lines 188--330)}\label{sec:fourier}
%======================================================================

\src{191--197}
\content Fourier conventions; Plancherel's identity; the inverse transform $F^\vee$; Hermitian functions.
\why We check the claims in order.
\begin{enumerate}
\item \emph{Plancherel.} For $g\in L^1\cap L^2$, $\|\hatf g\|_{L^2}^2=(2\pi)^d\|g\|_{L^2}^2$ with this
normalization. Since $L^1\cap L^2$ is dense in $L^2$, the map extends uniquely to $L^2$.
\item \emph{$F^\vee$ is continuous for $F\in L^1$.} The integrand $F(\xi)e^{i\xi\cdot x}$ is dominated by $|F(\xi)|$,
so continuity in $x$ follows by dominated convergence.
\item \emph{$F^\vee\in L^2$ and $\hatf{F^\vee}=F$ for $F\in L^1\cap L^2$.} We have $F^\vee(x)=(2\pi)^{-d}\hatf
F(-x)$. So $F^\vee\in L^2$ by Plancherel, and $\hatf{F^\vee}=F$ is the $L^2$ inversion theorem.
\item \emph{$g=(\hatf g)^\vee$ everywhere if $g\in L^2\cap C$ and $\hatf g\in L^1$.} Here $\hatf g\in L^1\cap L^2$,
so $(\hatf g)^\vee$ is continuous and has Fourier transform $\hatf g$. By injectivity of the $L^2$ Fourier transform,
$(\hatf g)^\vee=g$ almost everywhere. The set where two continuous functions differ is open; if it is also null, it
is empty.
\item \emph{$g$ real $\Rightarrow$ $\hatf g$ Hermitian.} For $g\in L^1$ we have
$\overline{\hatf g(\xi)}=\int g(x)e^{i\xi\cdot x}\dd x=\hatf g(-\xi)$. For $g\in L^2$, apply this to the truncations
$g\mathbf 1_{\{|x|\le n\}}$. Their transforms converge in $L^2$, hence almost everywhere along a subsequence, and the
identity passes to the limit.
\item \emph{$F\in L^1\cap L^2$ Hermitian $\Rightarrow$ $F^\vee$ real.} Using the substitution $\xi\to-\xi$,
\[
\overline{F^\vee(x)}=(2\pi)^{-d}\int\overline{F(\xi)}e^{-i\xi\cdot x}\dd\xi
=(2\pi)^{-d}\int F(-\xi)e^{-i\xi\cdot x}\dd\xi=F^\vee(x).
\]
\end{enumerate}
\verified \tA.

\src{198--201}
\content The Gaussian integral (2.1), for $P>0$:
\[
\int_{\R^d}e^{-x^\tp Px-i\xi\cdot x}\dd x=\pi^{d/2}(\det P)^{-1/2}e^{-\frac14\xi^\tp P^{-1}\xi}.
\]
\why \emph{One dimension.} Let $I(\xi):=\int_\R e^{-px^2-i\xi x}\dd x$ with $p>0$. Differentiating under the integral
(dominated by $|x|e^{-px^2}$) gives $I'(\xi)=\int(-ix)e^{-px^2-i\xi x}\dd x$. Since
$-ixe^{-px^2}=\frac{i}{2p}\frac{\mathrm d}{\mathrm dx}e^{-px^2}$, integration by parts gives
\[
I'(\xi)=\frac{i}{2p}\int\Bigl(\frac{\mathrm d}{\mathrm dx}e^{-px^2}\Bigr)e^{-i\xi x}\dd x
=-\frac{i}{2p}\int e^{-px^2}(-i\xi)e^{-i\xi x}\dd x=-\frac{\xi}{2p}I(\xi).
\]
Together with $I(0)=\sqrt{\pi/p}$ this gives $I(\xi)=\sqrt{\pi/p}\,e^{-\xi^2/(4p)}$.

\emph{Dimension $d$.} Write $P=Q\Lambda Q^\tp$ with $Q$ orthogonal and $\Lambda=\operatorname{diag}(\lambda_j)$,
$\lambda_j>0$. Substitute $x=Qu$ (so $|\det Q|=1$) and put $\zeta=Q^\tp\xi$. Then $x^\tp Px=\sum_j\lambda_ju_j^2$ and
$\xi\cdot x=\zeta\cdot u$, and the integral factorizes:
\begin{align*}
\prod_{j=1}^d\sqrt{\pi/\lambda_j}\,e^{-\zeta_j^2/(4\lambda_j)}
&=\pi^{d/2}(\det P)^{-1/2}e^{-\frac14\zeta^\tp\Lambda^{-1}\zeta}\\
&=\pi^{d/2}(\det P)^{-1/2}e^{-\frac14\xi^\tp P^{-1}\xi}.
\end{align*}
\verified \tA; \tS: S1 checks $d=1$.

\src{203--215}
\content The definitions (2.2) of $w_C$ and $\kappa_C$, the space $\FC{C}$, the transform of a kernel section (2.3),
and the integral (2.4).
\why \emph{(2.3).} Substitute $x=y+u$ and apply (2.1) with $P=C^{-2}$. Since $(\det C^{-2})^{-1/2}=\det C$ and
$P^{-1}=C^2$,
\begin{align*}
\hatf{k^C_y}(\xi)&=e^{-i\xi\cdot y}\int e^{-u^\tp C^{-2}u-i\xi\cdot u}\dd u\\
&=e^{-i\xi\cdot y}\pi^{d/2}\det C\,e^{-\frac14\xi^\tp C^2\xi}=\kappa_Ce^{-i\xi\cdot y}w_C(\xi)^{-1}.
\end{align*}
\emph{(2.4).} Apply (2.1) with $\xi=0$ and $P=\frac14C^2$. Then $\det P=4^{-d}(\det C)^2$, so
\[
\int w_C^{-1}\dd\xi=\pi^{d/2}\,2^d/\det C=(4\pi)^{d/2}/\det C .
\]
\verified \tA; \tS: S1 checks (2.3) and (2.4) for $d=1$.
\notes $\FC{C}$ is to be read modulo null sets, as usual for $L^2$-type spaces.

\src{217--230}
\content Lemma 2.1:
\begin{itemize}
\item[(a)] $\FC{C}\subset L^1\cap L^2$, with an explicit $L^1$ bound;
\item[(b)] $H_C$ consists of the continuous, real, $L^2$ functions with $\hatf f\in\FC{C}$, and its inner product is
(2.5);
\item[(c)] $f\mapsto\hatf f$ is a bijection from $H_C$ onto $\FC{C}$.
\end{itemize}

\src{232--238}
\content Proof of (a).
\why By Cauchy--Schwarz and (2.4),
\[
\|F\|_{L^1}=\int|F|w_C^{1/2}w_C^{-1/2}\le\Bigl(\int|F|^2w_C\Bigr)^{1/2}\Bigl(\frac{(4\pi)^{d/2}}{\det C}\Bigr)^{1/2}.
\]
Since $w_C\ge1$, also $\|F\|_{L^2}^2\le\int|F|^2w_C$.
\verified \tA.

\src{240--246}
\content $\mathcal H$ (the right-hand side of (b)) with the form (2.5) is a real inner product space.
\why \emph{Vector space}: $|\hatf f+\hatf g|^2\le2|\hatf f|^2+2|\hatf g|^2$.

\emph{Absolute convergence}:
\[
\int|\hatf f||\hatf g|w_C\le\Bigl(\int|\hatf f|^2w_C\Bigr)^{1/2}\Bigl(\int|\hatf g|^2w_C\Bigr)^{1/2}.
\]
\emph{Realness}: put $I(\xi)=\hatf f(\xi)\overline{\hatf g(\xi)}w_C(\xi)$. Hermitian symmetry and evenness of $w_C$
give $I(-\xi)=\overline{I(\xi)}$, so $\int I=\int I(-\xi)\dd\xi=\overline{\int I}$.

\emph{Definiteness}: $\ip ff=0$ forces $\hatf f=0$ almost everywhere, because $w_C>0$. Then $f=0$ in $L^2$, and $f=0$
everywhere by continuity.
\verified \tA.

\src{248--251}
\content The inversion formula (2.6) holds at every $x$.
\why $\hatf f\in L^1$ by (a); now apply item~4 of the conventions (lines 191--197).

\src{253--256}
\content $\mathcal H$ is complete.
\why Let $(f_n)$ be Cauchy in $\mathcal H$. By (2.5), $(\hatf{f_n})$ is Cauchy in $L^2(w_C\dd\xi)$, which is complete,
so $\hatf{f_n}\to F$ there.
\begin{itemize}
\item A subsequence converges almost everywhere. Since $\hatf{f_n}(-\xi)=\overline{\hatf{f_n}(\xi)}$ off a null set,
and the reflection of a null set is null, $F$ is Hermitian.
\item Thus $F\in\FC{C}\subset L^1\cap L^2$, so $f:=F^\vee$ is continuous, real, in $L^2$, and $\hatf f=F$.
\item Hence $f\in\mathcal H$, and $\|f_n-f\|_{\mathcal H}\to0$.
\end{itemize}
\verified \tA.

\src{258--260}
\content The kernel sections belong to $\mathcal H$.
\why By (2.3), $\hatf{k^C_y}$ is Hermitian:
$\hatf{k^C_y}(-\xi)=\kappa_Ce^{i\xi\cdot y}w_C(\xi)^{-1}=\overline{\hatf{k^C_y}(\xi)}$. Moreover
$\int|\hatf{k^C_y}|^2w_C=\kappa_C^2\int w_C^{-1}<\infty$, and $k^C_y$ is continuous, real and square integrable.
\notes The normalization is consistent:
\begin{align*}
\|k^C_y\|^2&=(2\pi)^{-d}\kappa_C^{-1}\kappa_C^2\int w_C^{-1}
=(2\pi)^{-d}\pi^{d/2}\det C\,\frac{(4\pi)^{d/2}}{\det C}\\
&=(2\pi)^{-d}(2\pi)^d=1=k^C(y,y).
\end{align*}
\tS: S1 checks this for $d=1$ and for general $d$.

\src{262--267}
\content The reproducing property.
\why By (2.3), $\overline{\hatf{k^C_y}(\xi)}=\kappa_Ce^{i\xi\cdot y}w_C(\xi)^{-1}$. So
\[
\ip f{k^C_y}_{\mathcal H}=(2\pi)^{-d}\kappa_C^{-1}\int\hatf f\,\kappa_Ce^{i\xi\cdot y}w_C^{-1}w_C\dd\xi
=(2\pi)^{-d}\int\hatf f(\xi)e^{i\xi\cdot y}\dd\xi=f(y),
\]
by (2.6).
\verified \tA.

\src{269--279}
\content $\mathcal H=H_C$ isometrically: the uniqueness part of the Moore--Aronszajn theorem.
\why The argument has four steps.
\begin{itemize}
\item \emph{Density of $H_0=\operatorname{span}\{k^C_y\}$.} In any RKHS, $f\perp H_0$ implies
$f(y)=\ip f{k^C_y}=0$ for all $y$, so $f=0$. This applies both to $\mathcal H$ and to $H_C$.
\item \emph{Equal norms on $H_0$.} By the reproducing property, in both spaces the squared norm of
$\sum_j\alpha_jk^C_{y_j}$ equals $\sum_{i,j}\alpha_i\alpha_jk^C(y_i,y_j)$.
\item \emph{Pointwise control.} $|h(x)|=|\ip h{k^C_x}|\le\|h\|\,\|k^C_x\|=\|h\|$, because $\|k^C_x\|^2=k^C(x,x)=1$.
\item \emph{Limits.} Let $f\in\mathcal H$ and choose $f_n\in H_0$ with $f_n\to f$ in $\mathcal H$. Then $(f_n)$ is
Cauchy in $H_C$ and converges there to some $g$. Pointwise limits coincide, so $g=f$ and the norms agree. The converse
inclusion is the same argument with the roles exchanged; it uses the completeness of $\mathcal H$ (lines 253--256).
\end{itemize}
Finally, polarization gives $\ip fg=\frac14(\|f+g\|^2-\|f-g\|^2)$ in the real space.
\verified \tA; \tL: the uniqueness of the RKHS with a given kernel is \cite{Aronszajn} and \cite{PaulsenRaghupathi}.

\src{281--284}
\content Proof of (c).
\why \emph{Injective}: $\hatf f$ determines $f$ in $L^2$, hence everywhere by continuity. \emph{Surjective}: for
$F\in\FC{C}$, the function $F^\vee$ is real, continuous and in $L^2$, with $\hatf{F^\vee}=F$, so $F^\vee\in H_C$ by
(b).

\src{286--288}
\content Lemma 2.1 is the Gaussian case of Wendland's Fourier description of native spaces \cite[Theorem
10.12]{Wendland}. A power-series description is in \cite{SHS}.
\why Wendland uses $\hatf f_W(\omega)=(2\pi)^{-d/2}\int f(x)e^{-i\omega\cdot x}\dd x$ and
$(f,g)=(2\pi)^{-d/2}\int\hatf f_W\overline{\hatf g_W}/\hatf\Phi_W$. Substituting $\hatf f_W=(2\pi)^{-d/2}\hatf f$ and
$\hatf\Phi_W=(2\pi)^{-d/2}\hatf\varphi$ gives
\[
(2\pi)^{-d/2}\,\frac{(2\pi)^{-d}\int\hatf f\,\overline{\hatf g}/\hatf\varphi}{(2\pi)^{-d/2}}
=(2\pi)^{-d}\int\hatf f\,\overline{\hatf g}/\hatf\varphi .
\]
For $\varphi=k^C(\cdot,0)$ we have $\hatf\varphi=\kappa_Cw_C^{-1}$, which gives (2.5).
\notes The theorem number and the normalization of \cite{Wendland} are quoted from memory of the book; the book was
not accessible here. The author should confirm them in the book. Lemma 2.1 is proved independently in the paper, so
nothing depends on the citation.

\src{290--297}
\content Lemma 2.2 (exponential multipliers). Let $S$ be symmetric. Then $\int|F|^2w_Ce^{\frac14\xi^\tp S\xi}<\infty$
for all $F\in\FC{C}$ if and only if $S\le0$. Otherwise a real, even witness supported in a double cone exists.

\src{299--316}
\content The proof.
\why If $S\le0$, the multiplier is at most $1$. Otherwise, let $\lambda>0$ be an eigenvalue with unit eigenvector $v$
and put $U=\{\theta\in\Sph^{d-1}:\theta^\tp S\theta>\lambda/2\}$.
\begin{itemize}
\item $U$ is relatively open, since $\theta\mapsto\theta^\tp S\theta$ is continuous. It is symmetric and contains
$v$, since $v^\tp Sv=\lambda$.
\item The cone $\Gamma=\{r\theta:r>0,\theta\in U\}$ is the preimage of $U$ under the continuous map
$\xi\mapsto\xi/|\xi|$ on $\R^d\setminus\{0\}$, hence open. On $\Gamma$ we have $\xi^\tp S\xi>\frac\lambda2|\xi|^2$.
\item $F_S=w_C^{-1/2}(1+|\xi|^2)^{-d/2}\mathbf 1_\Gamma$ is real and even, hence Hermitian. Moreover
$\int|F_S|^2w_C=\int_\Gamma(1+|\xi|^2)^{-d}<\infty$, because $r^{d-1}(1+r^2)^{-d}$ is bounded near $0$ and
$O(r^{-d-1})$ at $\infty$.
\end{itemize}
Multiplying by $e^{\frac14\xi^\tp S\xi}\ge e^{\lambda|\xi|^2/8}$ on $\Gamma$ (note $\frac14\cdot\frac\lambda2=\frac\lambda8$)
gives $|U|\int_0^\infty r^{d-1}(1+r^2)^{-d}e^{\lambda r^2/8}\dd r=\infty$. Here $|U|>0$ because $U$ is a nonempty open
subset of the sphere. For $d=1$, $U=\{-1,1\}$ with counting measure.
\verified \tA.
\notes The witness is chosen so that $\int|F_S|^2w_C$ converges only because of a polynomial factor. Any Gaussian
growth in a cone therefore destroys integrability. This is the precise sense of ``decays as slowly as membership
permits''.

\src{318--330}
\content Corollary 2.3: $H_{C_1}\subset H_{C_2}$ if and only if $C_1^2\ge C_2^2$, and then
$\|f\|_{C_2}\le(\det C_1/\det C_2)^{1/2}\|f\|_{C_1}$.
\why By Lemma 2.1(b),(c), inclusion means $\int|F|^2w_{C_2}<\infty$ for all $F\in\FC{C_1}$. Write
$w_{C_2}=w_{C_1}e^{\frac14\xi^\tp(C_2^2-C_1^2)\xi}$ and apply Lemma 2.2 with $S=C_2^2-C_1^2$. For the norm, use
$w_{C_2}\le w_{C_1}$ in (2.5):
\[
\|f\|_{C_2}^2\le\frac{\kappa_{C_1}}{\kappa_{C_2}}\|f\|_{C_1}^2=\frac{\det C_1}{\det C_2}\|f\|_{C_1}^2 .
\]
\verified \tA; \tL: \PS[Proposition 3.3] states ``$H_{C_1}\hookrightarrow H_{C_2}$ if and only if $C_1^2\ge C_2^2$. In
this case, we have $\|f\|_{C_2}\le(\det C_1/\det C_2)^{1/2}\|f\|_{C_1}$''. The statement and the constant are
identical.

%======================================================================
\section{Paper Section 3: The Koopman operators in Fourier variables (lines 333--409)}\label{sec:koop}
%======================================================================

\src{336--343}
\content The solution is $X_t=e^{At}x+Z_t$ with $Z_t\sim\mathcal N(0,\Sigma(t))$, and hence (3.1):
$(K^tf)(x)=\int f(e^{At}x+y)\gamma_t(\mathrm dy)$.
\why \emph{Variation of constants.} By It\^o's formula, $\mathrm d(e^{-At}X_t)=-Ae^{-At}X_t\dd t+e^{-At}\mathrm
dX_t=e^{-At}B\dd W_t$. So $X_t=e^{At}x+\int_0^te^{A(t-s)}B\dd W_s$.

\emph{Law of $Z_t$.} $Z_t$ is a Wiener integral of a deterministic integrand, hence centered Gaussian. By the It\^o
isometry and the substitution $r=t-s$,
\[
\mathbb E[Z_tZ_t^\tp]=\int_0^te^{A(t-s)}BB^\tp e^{A^\tp(t-s)}\dd s=\int_0^te^{Ar}BB^\tp e^{A^\tp r}\dd r=\Sigma(t).
\]
Formula (3.1) is the expectation of $f(e^{At}x+Z_t)$.
\verified \tA; \tL: \cite{Vatiwutipong} (distribution of the multivariate Ornstein--Uhlenbeck process). \tN: items 1
and 5.

\src{344--345}
\content Every $f\in H_C$ is continuous (Lemma 2.1(b)) and bounded by $\|f\|_C$.
\why $|f(x)|=|\ip f{k^C_x}|\le\|f\|_C\|k^C_x\|_C=\|f\|_C$.
\notes Correction C8 applies here.

\src{347--353}
\content Remark 3.1. The $L^2(\mu)$ reading of $K^tH_C\subset H_C$ in \PS\ agrees with the pointwise reading.
\why Under controllability, $\Sigma>0$, so $\mu=\mathcal N(0,\Sigma)$ has an everywhere positive density. Two
continuous functions that agree $\mu$-almost everywhere agree Lebesgue-almost everywhere, hence everywhere. The
$L^2(\mu)$ Koopman operator is defined by the same conditional expectation, so its class of $K^tf$ contains the
continuous function (3.1).
\verified \tL: \PS[Lemma 3.1] reads ``the RKHS $H_C$ is densely and continuously embedded in $L^2(\mu)$''. \PS[Section
2] reads ``$\Sigma$ is positive definite thanks to controllability of $(A,B)$''.

\src{355--362}
\content Lemma 3.2. For bounded, continuous $f\in L^2$, $K^tf$ is bounded, continuous and in $L^2$, and (3.2) holds:
\[
\hatf{K^tf}(\xi)=e^{-t\Tr A}\hatf f(e^{-A^\tp t}\xi)\exp\bigl(-\tfrac12\xi^\tp e^{-At}\Sigma(t)e^{-A^\tp t}\xi\bigr).
\]

\src{364--380}
\content The proof: a convolution step and a dilation step.
\why \emph{Continuity}: by dominated convergence in (3.1), with dominating constant $\sup|f|$.

\emph{Convolution.} $J_tg(z)=\int g(z+y)\gamma_t(\mathrm dy)$. Minkowski's integral inequality applied to $|g|$ gives
$\bigl\|\int|g(\cdot+y)|\gamma_t(\mathrm dy)\bigr\|_{L^2}\le\int\|g(\cdot+y)\|_{L^2}\gamma_t(\mathrm dy)=\|g\|_{L^2}$.
Hence the integral converges absolutely for almost every $z$, and $\|J_tg\|_{L^2}\le\|g\|_{L^2}$. For $g\in L^1\cap L^2$,
Fubini's theorem applies, because $\iint|g(z+y)|\dd z\,\gamma_t(\mathrm dy)=\|g\|_{L^1}$. With $u=z+y$,
\[
\int g(z+y)e^{-i\xi\cdot z}\dd z=e^{i\xi\cdot y}\hatf g(\xi),\qquad
\int e^{i\xi\cdot y}\gamma_t(\mathrm dy)=e^{-\frac12\xi^\tp\Sigma(t)\xi}.
\]
The second identity is the characteristic function of $\mathcal N(0,\Sigma(t))$; it holds also for singular
$\Sigma(t)$, which is how degenerate Gaussian measures are defined. Both sides of
$\hatf{J_tg}=\hatf g\,e^{-\frac12\xi^\tp\Sigma\xi}$ are bounded linear maps $L^2\to L^2$, since the multiplier is at
most $1$. So the identity extends by density.

\emph{Dilation.} For $u\in L^1\cap L^2$ and invertible $E$, substitute $z=Ex$, $\mathrm dx=|\det E|^{-1}\mathrm dz$,
and use $\xi\cdot E^{-1}z=(E^{-\tp}\xi)\cdot z$:
\[
\hatf{u\circ E}(\xi)=\int u(Ex)e^{-i\xi\cdot x}\dd x=|\det E|^{-1}\hatf u(E^{-\tp}\xi).
\]
Both sides are bounded on $L^2$: $\|u\circ E\|_{L^2}=|\det E|^{-1/2}\|u\|_{L^2}$. So the formula extends by density.

\emph{Assembly.} Take $E=e^{At}$. Jacobi's formula gives $\det e^{At}=e^{t\Tr A}>0$, and
$E^{-\tp}=e^{-A^\tp t}$. For bounded $f$, $J_tf$ is defined everywhere and agrees with the $L^2$ version almost
everywhere. Then $K^tf=(J_tf)\circ e^{At}$ yields (3.2).
\verified \tA; \tS: S2 checks the characteristic function and (3.2) for $f=k^C_z$, $d=1$. \tN: items 4, 5 and 9 use
consequences of (3.2).

\src{382--388}
\content The definitions (3.3) of $D_t$ and $\Lambda_C$, and three rewritings.
\why (1.4) is $D_t\ge0$ by definition. (1.5) is equivalent to $AC^2+C^2A^\tp\le4BB^\tp$, that is $\Lambda_C\ge0$.
Multiplying (1.3) by $2$ gives $2BB^\tp-(AC^2+C^2A^\tp)\ge0$.

\src{390--409}
\content Proposition 3.3, the norm identity (3.4), valid in $[0,\infty]$:
\[
\int|\hatf{K^tf}|^2w_C\dd\xi=e^{-t\Tr A}\int|\hatf f(\eta)|^2w_C(\eta)e^{-\frac14\eta^\tp D_t\eta}\dd\eta .
\]
\why Square (3.2) and multiply by $w_C(\xi)$ to get
\[
e^{-2t\Tr A}|\hatf f(e^{-A^\tp t}\xi)|^2\exp\bigl(-\xi^\tp e^{-At}\Sigma e^{-A^\tp t}\xi+\tfrac14\xi^\tp C^2\xi\bigr).
\]
Substitute
$\xi=e^{A^\tp t}\eta$, with $\mathrm d\xi=|\det e^{A^\tp t}|\dd\eta=e^{t\Tr A}\dd\eta$. The exponent becomes
\[
-\eta^\tp\Sigma(t)\eta+\tfrac14\eta^\tp e^{At}C^2e^{A^\tp t}\eta
=\tfrac14\eta^\tp C^2\eta-\tfrac14\eta^\tp\bigl(C^2+4\Sigma(t)-e^{At}C^2e^{A^\tp t}\bigr)\eta .
\]
The substitution is valid for nonnegative integrands with values in $[0,\infty]$.
\verified \tA; \tS: S2 checks that the integrands agree after the substitution ($d=1$, $f=k^C_z$). \tN: item 4 (the
value of $\|K^tk^C_z\|_C^2$).

%======================================================================
\section{Paper Section 4: Proofs of Theorem 1.1 and Corollary 1.2 (lines 412--585)}\label{sec:proofs}
%======================================================================

\src{415--426}
\content Theorem 4.1 (fixed time):
\begin{itemize}
\item[(a)] $K^tH_C\subset H_C\iff D_t\ge0$;
\item[(b)] if $D_t\ge0$, then $\|K^t\|=e^{-t\Tr A/2}$;
\item[(c)] if $D_t\ge0$ fails, a witness exists whose Fourier transform is real, even and supported in a cone.
\end{itemize}

\src{428--431}
\content $K^tf\in H_C$ if and only if the left side of (3.4) is finite.
\why $K^tf$ is real, continuous and in $L^2$ by Lemma 3.2 and (3.1). Its transform is Hermitian because it is real.
Now apply Lemma 2.1(b).

\src{433--434}
\content Sufficiency and the upper bound.
\why $D_t\ge0$ gives $e^{-\frac14\eta^\tp D_t\eta}\le1$. So (3.4) and (2.5) give
$\|K^tf\|_C^2\le e^{-t\Tr A}\|f\|_C^2$.

\src{436--444}
\content The lower bound, $\|K^t\|^2\ge e^{-t\Tr A}$.
\why $F_\varepsilon=\mathbf 1_{\{|\eta|<\varepsilon\}}$ is real and even, and $\int F_\varepsilon w_C<\infty$. So
$f_\varepsilon=F_\varepsilon^\vee\in H_C\setminus\{0\}$. On the ball, $\eta^\tp D_t\eta\le\|D_t\|\varepsilon^2$, so
\[
\|K^tf_\varepsilon\|_C^2=e^{-t\Tr A}(2\pi)^{-d}\kappa_C^{-1}\int_{|\eta|<\varepsilon}w_Ce^{-\frac14\eta^\tp D_t\eta}
\ge e^{-t\Tr A}e^{-\frac14\varepsilon^2\|D_t\|}\|f_\varepsilon\|_C^2 .
\]
Let $\varepsilon\downarrow0$.
\notes The norm is attained asymptotically by low-frequency functions. There $e^{-\frac14\eta^\tp D_t\eta}\approx1$,
and only the Jacobian factor $e^{-t\Tr A}$ of the dilation remains.

\src{446--450}
\content Necessity and the witness.
\why Apply Lemma 2.2 with $S=-D_t$. Then $S\le0$ fails, and the resulting $F_S$ gives
$\int|F_S|^2w_Ce^{-\frac14\eta^\tp D_t\eta}=\infty$. The function $f=F_S^\vee$ lies in $H_C$ with $\hatf f=F_S$, and
(3.4) shows that the left side is $\infty$, so $K^tf\notin H_C$.
\verified \tA.

\src{452--461}
\content Theorem 4.2: $K^t$ maps $H_C$ bijectively onto $H_{\Ct}$, and $\|K^tf\|_{\Ct}^2=\taut\|f\|_C^2$. In other
words, $\taut^{-1/2}K^t\colon H_C\to H_{\Ct}$ is unitary.
\why $\Ct$ is well defined and positive definite, because $C^2+4\Sigma(t)\ge C^2>0$ and congruence preserves
positivity.

\src{463--478}
\content Image and isometry.
\why By the computation of Proposition 3.3 with weight $w_{\Ct}$, the exponent is
$-\eta^\tp\Sigma\eta+\frac14\eta^\tp e^{At}\Ct^{\,2}e^{A^\tp t}\eta$. Since
$e^{At}\Ct^{\,2}e^{A^\tp t}=C^2+4\Sigma(t)$, it equals
\[
-\eta^\tp\Sigma\eta+\tfrac14\eta^\tp(C^2+4\Sigma)\eta=\tfrac14\eta^\tp C^2\eta .
\]
So $\int|\hatf g|^2w_{\Ct}=e^{-t\Tr A}\int|\hatf f|^2w_C$. Moreover $\det\Ct=e^{-t\Tr A}\det(C^2+4\Sigma)^{1/2}$,
hence
\[
\|g\|_{\Ct}^2=\frac{\kappa_C}{\kappa_{\Ct}}e^{-t\Tr A}\|f\|_C^2
=\frac{\det C\,e^{-t\Tr A}}{e^{-t\Tr A}\det(C^2+4\Sigma)^{1/2}}\|f\|_C^2=\taut\|f\|_C^2 .
\]
\verified \tA; \tS: S5 (the exponent identity, generic symmetric $2\times2$ matrices) and S2 (for $d=1$ and $f=k^C_z$,
$\|K^tk^C_z\|_{\Ct}^2=c/\sqrt{c^2+4\sigma}=\taut$, computed from an explicit Fourier integral). \tN: item 4 (Gram
matrices, to $3.8\times10^{-14}$).

\src{480--496}
\content Surjectivity.
\why Put $F(\eta)=e^{t\Tr A}\hatf g(e^{A^\tp t}\eta)e^{\frac12\eta^\tp\Sigma\eta}$. It is Hermitian: $\hatf g$ is
Hermitian and the other factors are real and even. With $\eta=e^{-A^\tp t}\xi$ and $\mathrm d\eta=e^{-t\Tr A}\mathrm
d\xi$,
\[
\eta^\tp\Sigma\eta+\tfrac14\eta^\tp C^2\eta=\tfrac14\eta^\tp(C^2+4\Sigma)\eta=\tfrac14\xi^\tp\Ct^{\,2}\xi .
\]
So $\int|F|^2w_C=e^{t\Tr A}\int|\hatf g|^2w_{\Ct}<\infty$. By (3.2),
\[
\hatf{K^tF^\vee}(\xi)=e^{-t\Tr A}\cdot e^{t\Tr A}\hatf g(\xi)e^{\frac12\xi^\tp e^{-At}\Sigma e^{-A^\tp t}\xi}
\cdot e^{-\frac12\xi^\tp e^{-At}\Sigma e^{-A^\tp t}\xi}=\hatf g(\xi),
\]
and continuity gives $K^tF^\vee=g$.
\verified \tA.

\src{498--509}
\content Corollary 4.3. $K^tH_C\subset H_C\iff H_{\Ct}\subset H_C\iff\Ct^{\,2}\ge C^2\iff D_t\ge0$. Moreover
$H_{\Ct}\subsetneq H_{C_t}$ when $\Sigma(t)\ne0$, and the implication discussed in \PS[Remark 4.2] fails for the matrices
of Corollary 1.2.

\src{511--516}
\content The chain of equivalences, and strictness.
\why $\Ct^{\,2}-C^2=e^{-At}(C^2+4\Sigma-e^{At}C^2e^{A^\tp t})e^{-A^\tp t}=e^{-At}D_te^{-A^\tp t}$, a congruence.
Next, $\Ct^{\,2}-C_t^2=e^{-At}(4\Sigma-2\Sigma)e^{-A^\tp t}=2e^{-At}\Sigma e^{-A^\tp t}\ge0$. If both inclusions held,
Corollary 2.3 would give $C_t^2\ge\Ct^{\,2}$ as well, so $\Sigma(t)\ge0$ and $-\Sigma(t)\ge0$, hence $\Sigma(t)=0$.
If $B\ne0$ and $t>0$, then $\Tr\Sigma(t)=\int_0^t\|B^\tp e^{A^\tp s}\|_F^2\dd s>0$, because the continuous integrand
equals $\|B\|_F^2>0$ at $s=0$. (Here $\|\cdot\|_F$ is the Frobenius norm.)
\verified \tA; \tS: S5 (the identity $\Ct^{\,2}-C_t^2=2e^{-At}\Sigma e^{-A^\tp t}$) and S3 (the congruence, Example
6.1).

\src{518--524}
\content The implication ``$K^tH_C\subset H_C\Rightarrow H_{C_t}\subset H_C$'' fails for the matrices of Corollary
1.2.
\why $f_x(t)=\ip{C_t^2x}x$ with $C_t^2=e^{-At}(C^2+2\Sigma(t))e^{-A^\tp t}$. Using $\Sigma(0)=0$ and
$\dot\Sigma(0)=BB^\tp$,
\begin{gather*}
\frac{\mathrm d}{\mathrm dt}C_t^2\Big|_{t=0}=-AC^2-C^2A^\tp+2BB^\tp,\\
\dot f_x(0)=\ip{(2BB^\tp-(AC^2+C^2A^\tp))x}x .
\end{gather*}
For $\beta=5$ the matrix $2I-(A+A^\tp)=\begin{pmatrix}4&-5\\-5&4\end{pmatrix}$ has the eigenvalue $-1$. So
$\dot f_x(0)<0$ for an eigenvector $x$, and $f_x(t)<f_x(0)=|Cx|^2$ on some $(0,\varepsilon)$. Corollary 2.3 then gives
$H_{C_t}\not\subset H_C$. Invariance for all $t$ comes from Theorem 4.4; there is no circularity, because Theorem 4.4
does not use Corollary 4.3.
\verified \tA; \tN: item 6. For $\beta\in\{5,6\}$, $C_t^2-C^2$ has the negative eigenvalues $-1.2\times10^{-2}$ and
$-2.4\times10^{-2}$ at some $t\le0.0124$. These match the first-order prediction $\lambda_{\min}\approx-t$ (resp.\
$-2t$) at $t=0.0124$.
\notes Correction C3 applies here.

\src{526--533}
\content Theorem 4.4. The following are equivalent: (i) invariance for all $t$; (ii) invariance along a sequence
$t_n\downarrow0$; (iii) $\Lambda_C\ge0$.

\src{535--549}
\content The proof via (4.2), $D_t=\int_0^te^{As}\Lambda_Ce^{A^\tp s}\dd s$.
\why Clearly $D_0=0$, because $\Sigma(0)=0$. By the fundamental theorem of calculus,
$\frac{\mathrm d}{\mathrm dt}4\Sigma(t)=4e^{At}BB^\tp e^{A^\tp t}$. By the product rule and $Ae^{At}=e^{At}A$,
$\frac{\mathrm d}{\mathrm dt}(e^{At}C^2e^{A^\tp t})=e^{At}(AC^2+C^2A^\tp)e^{A^\tp t}$. So
$\dot D_t=e^{At}\Lambda_Ce^{A^\tp t}$, and integrating gives (4.2).
\begin{itemize}
\item (iii)$\Rightarrow$(i): the integrand in (4.2) is PSD.
\item (ii)$\Rightarrow$(iii): $t_n^{-1}D_{t_n}\ge0$, and $t^{-1}D_t\to\Lambda_C$ by continuity of the integrand. The
PSD cone is closed.
\end{itemize}
\verified \tA; \tS: S3 checks (4.2), $\dot D_t$ and $\lim_{t\to0}D_t/t=\Lambda_C$ for Example 6.1. \tN: item 2, for
three random non-Hurwitz triples.

\src{551--556}
\content Remark 4.5. The same computation as in \PS, with $\Ct$ in place of $C_t$, gives
$\widetilde f_x(t)=\ip{\Ct^{\,2}x}x=|Ce^{-A^\tp t}x|^2+4\int_0^t|B^\tp e^{-A^\tp s}x|^2\dd s$ and
$\dot{\widetilde f}_x(t)=\ip{\Lambda_Ce^{-A^\tp t}x}{e^{-A^\tp t}x}$.
\why Put $y=e^{-A^\tp t}x$. Then $\ip{\Ct^{\,2}x}x=\ip{(C^2+4\Sigma(t))y}y$. For the $\Sigma$ term, substitute
$r=t-s$:
\[
\ip{\Sigma(t)y}y=\int_0^t|B^\tp e^{A^\tp s}e^{-A^\tp t}x|^2\dd s=\int_0^t|B^\tp e^{-A^\tp r}x|^2\dd r .
\]
Differentiating,
\begin{gather*}
\frac{\mathrm d}{\mathrm dt}|Ce^{-A^\tp t}x|^2=-\ip{(AC^2+C^2A^\tp)y}y,\\
\frac{\mathrm d}{\mathrm dt}\,4\int_0^t|B^\tp e^{-A^\tp r}x|^2\dd r=4\ip{BB^\tp y}y .
\end{gather*}
Together this is $\ip{\Lambda_Cy}y$.
\verified \tA; \tS: S3 checks $\frac{\mathrm d}{\mathrm dt}\Ct^{\,2}=e^{-At}\Lambda_Ce^{-A^\tp t}$ for Example 6.1.

\src{558--563}
\content The proof of Theorem 1.1.
\why (a) is Theorem 4.1. (b) is Theorem 4.4, (i)$\iff$(iii). For the interval statement: if every interval
$(0,1/n)$ contained an admissible $t_n$, then $t_n\to0$, and a strictly decreasing subsequence would satisfy (ii).

\src{565--572}
\content The proof of Corollary 1.2.
\why $BB^\tp\ge0$ gives $BB^\tp\le2BB^\tp$, so (1.3) implies (1.5). For the example, $A+A^\tp=\begin{pmatrix}-2&5\\5&-2
\end{pmatrix}$ has the eigenvalues $-2\pm5$, so $\frac12(A+A^\tp)$ has the eigenvalues $\frac32$ and $-\frac72$.
Hence (1.3) reads $\frac32\le1$, which is false, and (1.5) reads $\frac32\le2$, which is true. Moreover $A$ is Hurwitz
(double eigenvalue $-1$), $B=I$ is invertible (so $(A,B)$ is controllable), and $m=d=2$.
\verified \tA; \tS: S3.

\src{574--585}
\content Corollary 4.6, the isotropic case: invariance for all $t$ if and only if
$\frac12(A+A^\tp)\le\frac2{\sigma^2}BB^\tp$. This sharpens \PS[Corollary 2.3], which has $\frac1{\sigma^2}$.
\why Put $C=\sigma I$ in (1.5) and divide by $\sigma^2$.
\verified \tL: \PS[Corollary 2.3] reads ``assume that $\frac12(A+A^\top)\le\frac1{\sigma^2}BB^\top$ (e.g., if $A$ is
dissipative)''.

%======================================================================
\section{Paper Section 5: Stable drift (lines 588--642)}\label{sec:geom}
%======================================================================

\src{591--611}
\content Theorem 5.1, for Hurwitz $A$ and without controllability. Here $P_C=C^2+4\Sigma>0$, and
$\mathcal E_C=\{x:x^\tp P_C^{-1}x\le1\}$.
\begin{itemize}
\item[(a)] $D_t=P_C-e^{At}P_Ce^{A^\tp t}$, so invariance at time $t$ is equivalent to $e^{At}\mathcal E_C\subset
\mathcal E_C$.
\item[(b)] $AP_C+P_CA^\tp=-\Lambda_C$, so (1.5) is equivalent to forward invariance of $\mathcal E_C$; likewise for
(1.3) with $C^2+2\Sigma$.
\item[(c)] $T_C$ is a closed additive semigroup, contains $[t_C,\infty)$, and obeys the dichotomy.
\end{itemize}
\why $P_C\ge C^2>0$.

\src{613--619}
\content Proof of (a).
\why First,
\[
\int_t^\infty e^{As}BB^\tp e^{A^\tp s}\dd s=e^{At}\Bigl(\int_0^\infty e^{Ar}BB^\tp e^{A^\tp r}\dd r\Bigr)e^{A^\tp
t}=e^{At}\Sigma e^{A^\tp t},
\]
so $\Sigma(t)=\Sigma-e^{At}\Sigma e^{A^\tp t}$. Insert this into $D_t$:
\[
D_t=C^2+4\Sigma-4e^{At}\Sigma e^{A^\tp t}-e^{At}C^2e^{A^\tp t}=P_C-e^{At}P_Ce^{A^\tp t}.
\]
With $G=P^{-1/2}EP^{1/2}$ we have $GG^\tp=P^{-1/2}EPE^\tp P^{-1/2}$. Therefore
\[
EPE^\tp\le P\iff GG^\tp\le I\iff|G^\tp z|\le|z|\ \forall z\iff\|G\|=\|G^\tp\|\le1 .
\]
Finally, $|P^{-1/2}Ex|=|Ex|_P$ and $|P^{-1/2}x|=|x|_P$. By homogeneity, contraction in a norm is the same as mapping
its unit ball into itself.
\verified \tA; \tS: S3 (Example 6.1). \tN: item 3 (random Hurwitz $A$, entries of the difference at most
$2.4\times10^{-14}$).

\src{621--626}
\content Proof of (b).
\why $A\Sigma+\Sigma A^\tp=\int_0^\infty\frac{\mathrm d}{\mathrm ds}(e^{As}BB^\tp e^{A^\tp s})\dd s=0-BB^\tp$,
because $e^{As}\to0$. Hence $AP_C+P_CA^\tp=AC^2+C^2A^\tp-4BB^\tp=-\Lambda_C$, and similarly with $2\Sigma$. Forward
invariance means $e^{At}P_Ce^{A^\tp t}\le P_C$ for all $t$, by (a).
\begin{itemize}
\item If $AP_C+P_CA^\tp\le0$, then $\frac{\mathrm d}{\mathrm dt}e^{At}P_Ce^{A^\tp t}=e^{At}(AP_C+P_CA^\tp)e^{A^\tp
t}\le0$.
\item Conversely, divide $P_C-e^{At}P_Ce^{A^\tp t}\ge0$ by $t$ and let $t\downarrow0$.
\end{itemize}
\verified \tA; \tS: S3 and S4 ($AP_C+P_CA^\tp=-\Lambda_C$ exactly). \tN: item 3.

\src{628--633}
\content Proof of (c).
\why \emph{Closed and contains $0$}: $D_0=0$, $t\mapsto D_t$ is continuous, and the PSD cone is closed.

\emph{Additive}: congruence by $e^{As}$ preserves $\le$, so
$e^{A(s+t)}P_Ce^{A^\tp(s+t)}\le e^{As}P_Ce^{A^\tp s}\le P_C$.

\emph{Tail}: $x^\tp e^{At}P_Ce^{A^\tp t}x\le\lambda_{\max}(P_C)|e^{A^\tp t}x|^2\le\lambda_{\max}(P_C)\|e^{At}\|^2|x|^2$.
This is at most $\lambda_{\min}(P_C)|x|^2\le x^\tp P_Cx$ when $\|e^{At}\|^2\le\lambda_{\min}/\lambda_{\max}$. Such $t$
exist because $\|e^{At}\|\to0$.

\emph{Dichotomy}: Theorem 1.1(b).
\verified \tA.
\notes $T_C$ need not be an interval in general. The paper claims only the semigroup and tail structure, and the
interval shape in Example 6.1 is labelled as numerical.

\src{635--642}
\content Remark 5.2. The ellipsoid of $\Sigma$ is forward invariant whenever $\Sigma>0$. Invariance of $H_C$ is forward
invariance of the ellipsoid of $\Sigma+\frac14C^2$; condition (1.3) asks this for $\Sigma+\frac12C^2$. For Hurwitz $A$,
every $H_C$ is invariant for large $t$.
\why $\frac{\mathrm d}{\mathrm dt}e^{At}\Sigma e^{A^\tp t}=-e^{At}BB^\tp e^{A^\tp t}\le0$. The ellipsoids of $P$ and
of $cP$ ($c>0$) are dilates of each other, so forward invariance does not depend on $c$; here
$P_C=4(\Sigma+\frac14C^2)$ and $C^2+2\Sigma=2(\Sigma+\frac12C^2)$. The last sentence is Theorem 5.1(c).
\notes Correction C5 applies here.

%======================================================================
\section{Paper Section 6: Examples (lines 645--732)}\label{sec:ex}
%======================================================================

\src{648--661}
\content Example 6.1: $A_\beta=\begin{pmatrix}-1&\beta\\0&-1\end{pmatrix}$, $B=C=I_2$, $\beta\ge0$. Condition (1.3)
holds if and only if $\beta\le4$, and (1.5) if and only if $\beta\le6$. For $4<\beta\le6$ the space is invariant with
$\|K^t\|=e^t$; for $\beta>6$ invariance fails for small $t$.
\why $A+A^\tp=\begin{pmatrix}-2&\beta\\\beta&-2\end{pmatrix}$ has the eigenvalues $-2\pm\beta$. For $\beta\ge0$ the
largest eigenvalue of $\frac12(A+A^\tp)$ is $\frac{\beta-2}2$. So (1.3) reads $\frac{\beta-2}2\le1$, i.e.\
$\beta\le4$, and (1.5) reads $\frac{\beta-2}2\le2$, i.e.\ $\beta\le6$. Also $\Tr A_\beta=-2$ gives
$e^{-t\Tr A/2}=e^t$.
\verified \tA; \tS: S3 (eigenvalues, both thresholds, trace). \tN: item 6 ($\beta\in\{4,\dots,8\}$).

\src{663--668}
\content $e^{A_\beta t}$ and the integrand of $\Sigma(t)$.
\why $A_\beta=-I+\beta E_{12}$, and $-I$ commutes with $E_{12}$ while $E_{12}^2=0$. So
$e^{A_\beta t}=e^{-t}(I+\beta tE_{12})$, and
\[
e^{A_\beta s}e^{A_\beta^\tp s}=e^{-2s}\begin{pmatrix}1&\beta s\\0&1\end{pmatrix}
\begin{pmatrix}1&0\\\beta s&1\end{pmatrix}=e^{-2s}\begin{pmatrix}1+\beta^2s^2&\beta s\\\beta s&1\end{pmatrix}.
\]
\verified \tA; \tS: S3 (SymPy's matrix exponential agrees).

\src{669--678}
\content Three elementary integrals and the entries of $D_t$.
\why \emph{Antiderivatives.} The antiderivatives of $e^{-2s}$, $se^{-2s}$ and $s^2e^{-2s}$ are $-\frac12e^{-2s}$,
$-e^{-2s}(\frac s2+\frac14)$ and $-e^{-2s}(\frac{s^2}2+\frac s2+\frac14)$. For the last one, differentiate:
$e^{-2s}(s^2+s+\frac12)-e^{-2s}(s+\frac12)=s^2e^{-2s}$. Evaluating from $0$ to $t$ gives the three formulas.

\emph{The entries.} Write $e=e^{-2t}$. Then $4\Sigma_{11}=2(1-e)+\beta^2(1-e(1+2t+2t^2))$, and
$(e^{At}e^{A^\tp t})_{11}=e(1+\beta^2t^2)$. Hence
\begin{align*}
d_{11}&=1+2(1-e)+\beta^2\bigl(1-e(1+2t+2t^2)\bigr)-e(1+\beta^2t^2)\\
&=3(1-e)+\beta^2\bigl(1-e(1+2t+3t^2)\bigr).
\end{align*}
Similarly $d_{12}=\beta(1-e(1+2t))-e\beta t=\beta(1-e(1+3t))$ and $d_{22}=1+2(1-e)-e=3(1-e)$.
\verified \tA; \tS: S3 (the three integrals and the closed form of $D_t$, exactly). \tN: item 6 (closed form against
numerical $D_t$ to $2.5\times10^{-14}$).

\src{679--683}
\content $D_t=t\Lambda_C+O(t^2)$ with $\Lambda_C=\begin{pmatrix}6&-\beta\\-\beta&6\end{pmatrix}$.
\why $\Lambda_C=4I-(A+A^\tp)$. Also directly: $1-e^{-2t}=2t+O(t^2)$, $1-e^{-2t}(1+3t)=-t+O(t^2)$ and
$1-e^{-2t}(1+2t+3t^2)=O(t^2)$.
\verified \tS: S3 ($\lim_{t\to0}D_t/t=\Lambda_C$).

\src{684--694}
\content $\Sigma$, $P_C$, and the two Lyapunov matrices, which recover the thresholds $6$ and $4$.
\why Let $t\to\infty$ in the integrals: $\Sigma_{11}=\frac12+\frac{\beta^2}4$, $\Sigma_{12}=\frac\beta4$,
$\Sigma_{22}=\frac12$. So $P_C=I+4\Sigma=\begin{pmatrix}3+\beta^2&\beta\\\beta&3\end{pmatrix}$. Then
$A_\beta P_C=\begin{pmatrix}-3&2\beta\\-\beta&-3\end{pmatrix}$, and adding the transpose gives
$\begin{pmatrix}-6&\beta\\\beta&-6\end{pmatrix}$. This is negative semidefinite if and only if $\beta\le6$. Likewise
$A_\beta(I+2\Sigma)=\begin{pmatrix}-2&\frac32\beta\\-\frac12\beta&-2\end{pmatrix}$, and with its transpose
$\begin{pmatrix}-4&\beta\\\beta&-4\end{pmatrix}$, which is negative semidefinite if and only if $\beta\le4$.
\verified \tA; \tS: S3 (all four matrices; $A\Sigma+\Sigma A^\tp=-I$; Theorem 5.1(a) identity).

\src{696--703}
\content The numerically computed boundary $t_1(\beta)$ for $6<\beta\le14$, and a rigorous argument for $t\ge40$.
\why \emph{Numerical part.} For each of 320 values of $\beta$, \file{make\_figure.py} samples $\lambda_{\min}(D_t)$ on
$2\times10^5$ points of $(0,40]$. It asserts that the admissible grid points form a final segment, and bisects on
$\det D_t$. The results are $t_1(7)=0.300056$ and $t_1(8)=0.387611$.

\emph{Rigorous tail.} $\|e^{A_\beta t}\|\le e^{-t}(1+\beta t)$, and
$\frac{\mathrm d}{\mathrm dt}(1+\beta t)e^{-t}=e^{-t}(\beta-1-\beta t)<0$ for $t\ge1$. At $t=40$ and $\beta=14$,
$(1+560)e^{-40}=2.38\times10^{-15}<10^{-14}$. For $P_C=\begin{pmatrix}3+\beta^2&\beta\\\beta&3\end{pmatrix}$, the
ratio $r=\lambda_{\min}/\lambda_{\max}$ satisfies $r/(1+r)^2=\det P_C/(\Tr P_C)^2=(9+2\beta^2)/(6+\beta^2)^2$. This
function decreases in $\beta>0$: its derivative is $-4\beta(\beta^2+3)/(\beta^2+6)^3$. Since $r\mapsto r/(1+r)^2$
increases on $(0,1)$, $r$ also decreases in $\beta$. At $\beta=14$ the eigenvalues are $101\mp70\sqrt2$, and
\begin{align*}
\frac{101-70\sqrt2}{101+70\sqrt2}\ge\frac1{100}
&\iff9999\ge7070\sqrt2\\
&\iff9999^2=99\,980\,001\ge2\cdot7070^2=99\,969\,800,
\end{align*}
so $r\ge10^{-2}$ for $\beta\le14$. Then $\lambda_{\max}\|e^{At}\|^2<\lambda_{\min}$ for $t\ge40$, and the proof of
Theorem 5.1(c) gives $D_t>0$ strictly.
\verified \tS: S3 checks the monotonicity, the exact eigenvalues, the inequality $9999^2\ge2\cdot7070^2$ and the bound
$561e^{-40}<10^{-14}$ with 50 digits. \tN: item 7 and \file{figure\_output.txt}.
\notes The ratio at $\beta=14$ is $0.010026$, so the margin over $10^{-2}$ is small but exact. The margin in the other
factor is enormous ($10^{-28}$ against $10^{-2}$). The conclusion ``$T_I=\{0\}\cup[t_1(\beta),\infty)$'' is stated
``up to the grid resolution'', which is accurate.

\src{706--716}
\content Figure 1: the admissible region in the $(\beta,t)$ plane. The region $\beta\le6$ is rigorous (Theorem 4.4),
the hatched region $\beta\le4$ is covered by (1.3), and the curve $t_1(\beta)$ for $\beta>6$ is numerical.
\verified \tN: \file{make\_figure.py}; the legend references ``(1.3)'' and ``Theorem 4.1'' match the compiled numbering.
The fonts are embedded as TrueType, not Type~3.

\src{718--732}
\content Example 6.2, scalar noise: $A=\begin{pmatrix}\frac32&-2\\2&-2\end{pmatrix}$, $B=e_1$, $C=I$. Here (1.3)
fails and (1.5) holds.
\why $\Tr A=-\frac12$ and $\det A=-3+4=1$, so $\lambda^2+\frac12\lambda+1=0$, i.e.\
$\lambda=-\frac14\pm\frac{\sqrt{15}}4i$. Next, $AB=(\frac32,2)^\tp$ and $\det[B,AB]=\det\begin{pmatrix}1&\frac32\\0&2
\end{pmatrix}=2$. Further $\frac12(A+A^\tp)=\operatorname{diag}(\frac32,-2)$ and $BB^\tp=\operatorname{diag}(1,0)$.
\notes The paper does not record the following, which gives an independent confirmation through Theorem 5.1(b). The
stationary covariance is $\Sigma=\begin{pmatrix}5&4\\4&4\end{pmatrix}>0$. Indeed
$A\Sigma=\begin{pmatrix}-\frac12&-2\\2&0\end{pmatrix}$, so $A\Sigma+\Sigma A^\tp=\operatorname{diag}(-1,0)=-BB^\tp$. Then
$P_C=\begin{pmatrix}21&16\\16&17\end{pmatrix}$ and $AP_C+P_CA^\tp=\operatorname{diag}(-1,-4)\le0$.
\verified \tA; \tS: S4 (all of the above, exactly). \tN: item 8 ($\lambda_{\min}(D_t)\ge0$ on 3000 grid points).

%======================================================================
\section{Paper Section 7: Further remarks (lines 735--863)}\label{sec:remarks}
%======================================================================

\src{738--755}
\content Lemma 7.1. $k^{C'}_y\in H_{C''}$ if and only if $G=2C'^2-C''^2>0$, and then
\[
\ip{k^{C'}_y}{k^{C'}_{y'}}_{C''}=\frac{(\det C')^2}{\det C''(\det G)^{1/2}}e^{-(y-y')^\tp G^{-1}(y-y')} .
\]
\why By (2.3), $|\hatf{k^{C'}_y}|^2w_{C''}=\kappa_{C'}^2e^{-\frac12\xi^\tp C'^2\xi+\frac14\xi^\tp C''^2\xi}
=\kappa_{C'}^2e^{-\frac14\xi^\tp G\xi}$. Diagonalizing $G$ shows that $\int e^{-\frac14\xi^\tp G\xi}\dd\xi$ factorizes
into one-dimensional integrals, which are all finite if and only if every eigenvalue of $G$ is positive. Next, (2.1)
with $P=\frac14G$ gives $\pi^{d/2}2^d(\det G)^{-1/2}e^{-(y-y')^\tp G^{-1}(y-y')}$. The constants combine to
\[
(2\pi)^{-d}\frac{\pi^d(\det C')^2}{\pi^{d/2}\det C''}\pi^{d/2}2^d(\det G)^{-1/2}
=\frac{(\det C')^2}{\det C''(\det G)^{1/2}} .
\]
For $C'=C''=C$ we get $G=C^2$ and recover $k^C(y,y')$, a consistency check.
\verified \tA; \tS: S1 (the constant for general $d$; the full formula for $d=1$).

\src{757--772}
\content Remark 7.2.
\begin{itemize}
\item The kernel-section formula, derived from (3.2) and (2.3).
\item $K^tk^C_z\in H_C\iff2C_t^2-C^2>0$.
\item For small $t$, finite kernel expansions are mapped into $H_C$ whatever the sign of $D_t$.
\end{itemize}
\why The derivation is the one given for lines 148--160 in Section~\ref{sec:intro}. Lemma 7.1 with $C'=C_t$ and
$C''=C$ gives the membership criterion. Next,
\[
2C_t^2-C^2=e^{-At}(2C^2+4\Sigma-e^{At}C^2e^{A^\tp t})e^{-A^\tp t}=e^{-At}(C^2+D_t)e^{-A^\tp t}.
\]
Since $D_t\to0$ and $C^2>0$, this is positive definite for small $t$.

For a finite expansion $f=\sum_j\alpha_jk^C_{y_j}$, (2.3) gives $|\hatf f|\le\kappa_C(\sum_j|\alpha_j|)w_C^{-1}$. So
$|\hatf f|^2w_C\le c\,w_C^{-1}$, and
\[
e^{-t\Tr A}\int|\hatf f|^2w_Ce^{-\frac14\eta^\tp D_t\eta}\le c\,e^{-t\Tr A}\int e^{-\frac14\eta^\tp(C^2+D_t)\eta}<\infty
\]
whenever $C^2+D_t>0$.
\verified \tA; \tL: \PS[Remark 4.2] reads ``Note that $K^tH_C\subset H_C$ at least implies
$H_{0,C_t}=\operatorname{span}\{k^{C_t}_x:x\in\R^d\}\subset H_C$, see Proposition 4.1.''
\notes Corrections C2 and C4 apply here.

\src{774--791}
\content Remark 7.3.
\begin{itemize}
\item Theorem 4.2 refines \PS[Proposition 4.1]: it recovers the bound $\|K^tf\|_{C_t}^2\le\tau_t\|f\|_C^2$.
\item On kernel sections, $\ip{K^tk^C_z}{K^tk^C_{z'}}_{\Ct}=\taut k^C(z,z')$.
\end{itemize}
\why \emph{The bound.} Corollary 2.3 with $C_1=\Ct$ and $C_2=C_t$ gives the factor $\det\Ct/\det C_t$. Then
\[
\frac{\det\Ct}{\det C_t}\taut=\frac{e^{-t\Tr A}\det(C^2+4\Sigma)^{1/2}}{e^{-t\Tr A}\det(C^2+2\Sigma)^{1/2}}\cdot
\frac{\det C}{\det(C^2+4\Sigma)^{1/2}}=\tau_t .
\]
\emph{The kernel identity.} Apply Lemma 7.1 with $C'=C_t$, $C''=\Ct$ and $G=2C_t^2-\Ct^{\,2}=e^{-At}C^2e^{-A^\tp t}$.
Then $G^{-1}=e^{A^\tp t}C^{-2}e^{At}$, and the exponent becomes $(z-z')^\tp C^{-2}(z-z')$. Let
$a=\det(C^2+2\Sigma)$, $b=\det(C^2+4\Sigma)$, $c=\det C$ and $e=e^{t\Tr A}$. Then $\tau_t^2=c^2/a$,
$(\det C_t)^2=a/e^2$, $\det\Ct=\sqrt b/e$ and $(\det G)^{1/2}=c/e$. So the constant is
\[
\frac{c^2}{a}\cdot\frac{a/e^2}{(\sqrt b/e)(c/e)}=\frac c{\sqrt b}=\taut .
\]
\verified \tA; \tS: S5 checks both determinant identities and the formula for $2C_t^2-\Ct^{\,2}$. \tN: item 4
(Gram matrices). \tL: the statement of \PS[Proposition 4.1] gives the norm bound $\tau_t^{1/2}$.

\src{793--800}
\content Remark 7.4, for $B=0$. The composition operator $f\mapsto f\circ E$ maps $H_C$ into itself if and only if
$EC^2E^\tp\le C^2$, i.e.\ $\|C^{-1}EC\|\le1$, and then its norm is $|\det E|^{-1/2}$. The remark also cites
\cite{IIS} and the Fock-space analogue \cite{CarswellMacCluerSchuster}.
\why The dilation formula of Lemma 3.2 gives $\int|\hatf{f\circ E}|^2w_C=|\det E|^{-1}\int|\hatf f|^2w_C
e^{-\frac14\eta^\tp(C^2-EC^2E^\tp)\eta}$. From here the proof of Theorem 4.1 runs verbatim, with
$D=C^2-EC^2E^\tp$. Moreover $C^{-1}EC(C^{-1}EC)^\tp=C^{-1}EC^2E^\tp C^{-1}$, and
$\|M\|\le1\iff MM^\tp\le I$.
\verified \tA; \tL: the abstract of \cite{IIS} (arXiv:1911.11992) says that only affine maps induce bounded
composition operators on a large class of RKHSs of analytic positive definite functions. The Fock-space statement is
the main theorem of \cite{CarswellMacCluerSchuster}; its bibliographic data were checked earlier (Section~\ref{sec:bib}).
\notes Translations $f\mapsto f(\cdot+b)$ are isometries of $H_C$, because they multiply $\hatf f$ by $e^{i\xi\cdot b}$.
So every affine map $Ex+b$ with $\|C^{-1}EC\|\le1$ gives a bounded operator. This could be added in one sentence.

\src{802--807}
\content The Fourier description of $H_\varphi$ for a translation-invariant kernel, in the paper's normalization.
\why This is the normalization conversion given for lines 286--288 in Section~\ref{sec:fourier}.

\src{809--820}
\content Proposition 7.5. $K^tH_\varphi\subset H_\varphi$ if and only if $M_t=\sup m_t<\infty$, and then
$\|K^t\|^2=e^{-t\Tr A}M_t$. For the Gaussian kernel, $m_t=e^{-\frac14\eta^\tp D_t\eta}$.
\why \emph{Gaussian case}: $\hatf\varphi=\kappa_Cw_C^{-1}$, so
\[
m_t(\eta)=\exp\bigl(-\tfrac14\eta^\tp C^2\eta-\eta^\tp\Sigma\eta+\tfrac14\eta^\tp e^{At}C^2e^{A^\tp t}\eta\bigr)
=e^{-\frac14\eta^\tp D_t\eta}.
\]
\verified \tS: S2 ($d=1$). \tN: item 9 ($1000$ random $\eta$, relative error $3.2\times10^{-14}$).

\src{822--832}
\content The norm formula.
\why The substitution of Proposition 3.3 with weight $1/\hatf\varphi$ gives
\[
\|K^tf\|^2=(2\pi)^{-d}e^{-t\Tr A}\int|\hatf f|^2m_t/\hatf\varphi .
\]
\begin{itemize}
\item \emph{Upper bound}: $m_t\le M_t$.
\item \emph{Lower bound}: $\varphi$ is real and positive definite, hence even, so $\hatf\varphi$ and $m_t$ are even; and
$m_t$ is continuous, since $\hatf\varphi$ is continuous and positive. Choose $\eta_0$ with $m_t(\eta_0)>M_t-\delta$,
where $0<\delta<M_t$. For small $r$, $m_t>M_t-\delta$ on $U_r=B_r(\eta_0)\cup B_r(-\eta_0)$. The function
$f_r=(\mathbf 1_{U_r})^\vee$ lies in $H_\varphi$, because $\mathbf 1_{U_r}$ is real and even and $1/\hatf\varphi$ is
bounded on the compact set $\overline{U_r}$. It satisfies the stated inequality.
\end{itemize}
\verified \tA.

\src{834--842}
\content If $M_t=\infty$, a witness exists.
\why The sets $\{m_t>16^k\}$ are open, nonempty and symmetric. Choose bounded symmetric $E_k$ of positive measure in
them (for example two small balls), and put $a_k=(\int_{E_k}\hatf\varphi^{-1})^{-1/2}$. Then:
\begin{itemize}
\item each term of $F$ has norm $2^{-k}$ in $L^2(\hatf\varphi^{-1}\dd\xi)$, so by Minkowski $\|F\|\le\sum_{k\ge1}2^{-k}=1$;
\item $F\in L^1$ by Cauchy--Schwarz with $\int\hatf\varphi<\infty$;
\item $F\in L^2$ because $0<\hatf\varphi\le\|\varphi\|_{L^1}$;
\item on $E_k$, $m_t>16^k$, so $\int F^2m_t/\hatf\varphi\ge4^{-k}a_k^216^k\int_{E_k}\hatf\varphi^{-1}=4^k$ for every
$k$.
\end{itemize}
\verified \tA.

\src{844--856}
\content Example 7.6. For $\hatf\varphi=(1+|\xi|^2)^{-s}$ with $s>d/2$ (Mat\'ern/Sobolev), every $K^t$ leaves
$H^s$ invariant, with $\|K^t\|\le e^{-t\Tr A/2}\max\{1,\|e^{At}\|\}^s$.
\why $\hatf\varphi$ is integrable if and only if $2s>d$. Next,
$(1+a|\eta|^2)/(1+|\eta|^2)\le\max\{1,a\}$ with $a=\|e^{A^\tp t}\|^2=\|e^{At}\|^2$. With
$e^{-\eta^\tp\Sigma\eta}\le1$ this gives $m_t\le\max\{1,\|e^{At}\|^2\}^s$, and the square root gives the norm
bound.
\verified \tA; \tN: item 10 ($8.04\times10^4$ points, non-Hurwitz $A$).
\notes Correction C6 applies here. For kernels whose spectral density decays like $e^{-|\xi|}$, the Gaussian factor
$e^{-\eta^\tp\Sigma(t)\eta}$ dominates once $\Sigma(t)>0$. So the Lyapunov-type restriction is specific to spectral
densities with (at least) Gaussian decay.

\src{858--863}
\content Remark 7.7 (Outlook): the method needs the ``dilation times Gaussian'' structure, which fails for nonlinear
drifts or state-dependent noise.

%======================================================================
\section{Appendix A, the AI-use statement and data availability (lines 865--939)}\label{sec:app}
%======================================================================

\src{871--876}
\content The three scripts. Errors are measured in $|M|_{\max}$. $\Sigma(t)$ is computed by Van Loan's formula on steps
of length at most $0.25$, combined with the recursion $\Sigma(s+h)=\Sigma(h)+e^{Ah}\Sigma(s)e^{A^\tp h}$.
\why The recursion is $\Sigma(s+h)=\int_0^h+\int_h^{s+h}$, and the substitution $u=r-h$ in the second integral
gives $e^{Ah}\Sigma(s)e^{A^\tp h}$. The one-step formula multiplies $e^{-At}$ by $e^{At}$ and loses all digits for
large $t$; this was observed during development and fixed.
\verified \tN: item 1.

\src{877--925}
\content Items 1--10 and totals. Each number in the paper was compared with the script output of 24 September 2026.

\medskip
{\small
\begin{longtable}{L{26pt}L{136pt}L{118pt}L{32pt}}
\toprule
Item & Claim in the paper & Script output & Match\\
\midrule
\endhead
1 & $|\Sigma-\Sigma_{\mathrm q}|_{\max}\le2.7\times10^{-13}|\Sigma_{\mathrm q}|_{\max}$; $3\times3$, $3\times2$ &
max $2.656\times10^{-13}$ & yes\\
2 & $\le5.2\times10^{-13}\max\{1,|R_t|_{\max}\}$; abscissae $1.645$, $1.556$, $2.043$; $3.7\times10^{-6}$ at
$t=10^{-6}$ & max $5.177\times10^{-13}$; same abscissae; max $3.638\times10^{-6}$ & yes\\
3 & entries $\le2.4\times10^{-14}$ and $\le1.8\times10^{-14}$ & $2.309\times10^{-14}$; $1.776\times10^{-14}$ & yes\\
4 & Gram entries within $3.8\times10^{-14}$; relative error $1.1\times10^{-13}$ & $3.747\times10^{-14}$;
$1.078\times10^{-13}$ & yes\\
5 & $0.436373\pm2.2\times10^{-4}$ vs.\ $0.436051$; $1.45$ standard errors & same; ratio to 5 standard errors $0.2901$ &
yes\\
6 & (1.3) only for $\beta=4$; (1.5) for $\beta\le6$; indefinite segments end at $0.298$, $0.385$; $-1.2\times10^{-2}$,
$-2.4\times10^{-2}$; $2.5\times10^{-14}$ & as claimed; $0.2980$, $0.3847$; $-1.215\times10^{-2}$,
$-2.415\times10^{-2}$; $2.487\times10^{-14}$ & yes\\
7 & 320 values; interval; $t_1$ from $0.053$ to $0.537$ & $0.0526$; $0.5366$; assertions pass & yes\\
8 & eigenvalues $-0.25\pm0.968i$; $\lambda_{\max}=\pm0.5$; 3000 points & same; $\min\lambda_{\min}=10^{-4}$ & yes\\
9 & relative error $3.2\times10^{-14}$ & $3.197\times10^{-14}$ & yes\\
10 & $8.04\times10^4$ points, $10^{-3}\le|\eta|\le10^6$, abscissa $0.857$ & 402 radii $\times$ 200 unit
directions; same abscissa & yes\\
Total & 38 checks; 52 exact checks & 38/38; 52/52 & yes\\
\bottomrule
\end{longtable}
}
\notes Correction C10 applies here. Before the audit, item 2 wrote the spectral norm although the script uses the
entrywise maximum. The parameters (dimensions, grids, seeds) were read off the code of \file{verify.py}. Its random
generator is \texttt{numpy.random.default\_rng(20260924)}.

\src{927--935}
\content A commented AI-use statement, to be activated and adapted by the author.
\notes This is a blocking item; see Section~\ref{sec:ready}. The statement names the tool (Claude, Anthropic), what it
was used for, and the author's responsibility. Its last sentence must be made true by the author before it is
uncommented.

\src{937--939}
\content Data availability: no external data; the scripts are supplementary material.

%======================================================================
\section{Bibliography (lines 941--1019)}\label{sec:bib}
%======================================================================

For each entry: what it supports, and how the bibliographic data were checked. ``Web'' means a web search in this or
an earlier step of the session. ``PDF'' means the arXiv PDF of \PS\ supplied with the task.

\begin{itemize}
\item \cite{Aronszajn}: Aronszajn 1950 supports the Moore--Aronszajn uniqueness argument (lines 269--279).
Bibliographic data: standard; web.
\item \cite{CarswellMacCluerSchuster}: Carswell--MacCluer--Schuster, Fock-space composition operators (Remark 7.4).
Web.
\item \cite{Hertel}: Hertel--Philipp--Schaller--Worthmann (new) supports lines 77 and 855. Web: arXiv:2512.20247,
DOI 10.1137/26M1839689 \cite{src-hertel}. \emph{Caveat}: volume 25, issue 3 and pages 2044--2083 come from search
metadata, because the SIAM page is blocked here. Confirm them at the DOI before submission.
\item \cite{IIS}: Ikeda--Ishikawa--Sawano (new), Remark 7.4. Web: \emph{J.\ Math.\ Anal.\ Appl.} 511 (2022), no.~1,
126048; arXiv:1911.11992 \cite{src-iis}.
\item \cite{Klus}: Klus--Schuster--Muandet, \emph{J.\ Nonlinear Sci.} 30 (2020). Web.
\item \cite{Koehne}: Köhne et al.\ (updated) supports lines 77 and 855. Web: DOI 10.1137/24M1650120, SIADS 24 (2025)
501--529 \cite{src-koehne}. \emph{Caveat}: the page range comes from search metadata; confirm it at the DOI.
\item \cite{Koopman}: Koopman 1931. Web.
\item \cite{Kostic}: Kostic et al., NeurIPS 35 (2022). Web.
\item \cite{Mauroy}: Mauroy--Mezi\'c--Susuki (eds.), LNCIS vol.\ 484, Springer 2020. Web \cite{src-mauroy}.
\item \cite{PaulsenRaghupathi}: Paulsen--Raghupathi 2016, Cambridge Studies in Advanced Mathematics 152. Web.
\item \cite{PhilippACHA}: Philipp et al., \emph{Appl.\ Comput.\ Harmon.\ Anal.} 71 (2024) 101657. Web.
\item \cite{PSWPN}: \PS, PAMM 25 (2025) e202400127 and arXiv:2405.14429. PDF for all numbered statements; web for the
PAMM data \cite{src-pswpn}. The entry now notes that numbered statements follow arXiv v1. The PAMM version could not be
opened (Wiley is blocked here), so its numbering may differ.
\item \cite{SHS}: Steinwart--Hush--Scovel, \emph{IEEE Trans.\ Inform.\ Theory} 52 (2006). Web.
\item \cite{VanLoan}: Van Loan, \emph{IEEE Trans.\ Automat.\ Control} 23 (1978). Web. This is the block-exponential
method used in the scripts.
\item \cite{Vatiwutipong}: Vatiwutipong--Phewchean, \emph{Adv.\ Difference Equ.} 2019, Paper No.\ 276. Web.
\item \cite{Wendland}: Wendland 2005, Cambridge Monographs 17. Web for the book. The content of Theorem 10.12 is from
memory; see the note on lines 286--288.
\item \cite{WilliamsEDMD}, \cite{WilliamsKernel}: Williams--Kevrekidis--Rowley, \emph{J.\ Nonlinear Sci.} 25 (2015);
Williams--Rowley--Kevrekidis, \emph{J.\ Comput.\ Dyn.} 2 (2015). Web.
\end{itemize}
Every entry is cited in the text, and every citation has an entry (\tB: no undefined citations in the log). The keys
are ordered so that citation lists appear in increasing numerical order.

%======================================================================
\section{Questions a referee may ask}\label{sec:referee}
%======================================================================

\begin{enumerate}
\item \emph{Is $K^t$ well defined on $H_C$, and is it the operator of \PS?} Yes. $K^tf$ is defined pointwise by (3.1),
and Remark 3.1 shows agreement with the $L^2(\mu)$ operator of \PS\ under their assumptions.
\item \emph{Is the proof not just a Fourier computation?} The proof is short. Its value is threefold:
\begin{itemize}
\item it settles a published question with a sharp answer (Theorem 1.1);
\item it identifies the range of $K^t$ exactly (Theorem 4.2);
\item it pinpoints where the earlier argument loses information (Corollary 4.3, Remarks 7.2 and 7.3).
\end{itemize}
A cover letter should say this plainly.
\item \emph{What is the relation to the literature on Ornstein--Uhlenbeck semigroups?} The paper does not discuss it.
A referee from probability may ask for a paragraph relating (3.2) to the Mehler formula and to the classical theory of
Ornstein--Uhlenbeck semigroups on $L^p(\mu)$. Such a paragraph would need references that were not verified in this
audit, so none were added.
\item \emph{Do the standing assumptions matter?} No. Theorem 1.1 holds for all $A$, $B$ and $m$. Corollary 1.2 and
Examples 6.1 and 6.2 satisfy the standing assumptions, so the answer to \PS\ is given within their setting.
\item \emph{Is $T_C$ always an interval $\{0\}\cup[t_1,\infty)$?} The paper does not claim this. It proves the
semigroup and tail structure (Theorem 5.1(c)) and reports the interval shape in Example 6.1 as numerical.
\item \emph{What about nonlinear SDEs?} This is outside the scope, and Remark 7.7 says so.
\item \emph{How does this relate to \cite{Hertel}?} That paper treats Sobolev spaces for stochastic dynamics. Example
7.6 is the elementary linear case, and the Gaussian question is not addressed there.
\end{enumerate}

%======================================================================
\section{Submission readiness and choice of journal}\label{sec:ready}
%======================================================================

\subsection{Verdict}
The \emph{content} is ready:
\begin{itemize}
\item the mathematics is complete and every step was re-derived;
\item the claims are checked symbolically and numerically;
\item the literature is up to date as far as the searches reach;
\item the manuscript compiles without warnings.
\end{itemize}
The \emph{submission} is not ready, because of the blocking items below. None of them is mathematical.

\subsection{Blocking items}
\begin{enumerate}
\item \emph{Authorship.} Fill in the author block (lines 38--41): name, affiliation, e-mail and ORCID. The submitting
author must be able to defend every line. SIAM's editorial policy (version 2.0, May 2026) states that authors must be
human and able to directly defend all aspects of the submission \cite{src-siam-ai}. This commentary is meant to make
that possible. If the author cannot follow a derivation here, a colleague in functional analysis or kernel methods
should read the paper before submission.
\item \emph{AI-use disclosure.} Activate and adapt lines 927--935 according to the journal's policy.
\begin{itemize}
\item SIAM requires disclosure of substantive AI use, and purposeful non-disclosure is treated as misconduct
\cite{src-siam-ai,src-siam-news}.
\item Elsevier requires a separate AI declaration before the references \cite{src-elsevier}.
\item For the AMS I could not find the current policy; check it before submitting there.
\end{itemize}
\item \emph{Journal format and declarations.} Use the journal's class: \texttt{siamart} for SIAM, \texttt{elsarticle}
for Elsevier, the AMS \texttt{proc} class for PAMS. Add funding, competing-interest and code-availability statements as
required.
\item \emph{Cover letter.} It should name the open question, its source (\PS[Remark 2.2(c), Remark 4.2]), the answer,
and the three contributions listed in Section~\ref{sec:referee}, item 2.
\end{enumerate}

\subsection{Recommended items}
\begin{enumerate}
\item Post the paper on arXiv (math.FA, cross-listed to math.PR and math.DS) before or at submission, to establish
priority.
\item Repeat the novelty search on arXiv and MathSciNet/zbMATH, which could not be reached from here
(Section~\ref{sec:novelty}).
\item Confirm the volume and pages of \cite{Hertel} and \cite{Koehne} at their DOIs, and Theorem 10.12 of
\cite{Wendland} in the book.
\item Consider informing Philipp et al.\ as a courtesy. They posed the question and are natural referees.
\item Supply the three scripts and their outputs as supplementary material; the README lists them.
\item Optional additions: MSC 47B32; the observation that (1.5) is the condition of \PS\ for the squared kernel
(Section~\ref{sec:intro}); the remark on affine maps (Section~\ref{sec:remarks}).
\end{enumerate}

\subsection{Choice of journal}
\emph{First choice: SIAM Journal on Applied Dynamical Systems (SIADS).}
\begin{itemize}
\item Its scope includes contributions ``to the theory of dynamical systems'' \cite{src-siads}.
\item The error analysis of kernel EDMD that motivated the question is published there: Köhne et al.\ (2025) and
Hertel et al.\ (2026). The readers who need the result read SIADS.
\item SIAM's policy permits AI-assisted research with disclosure and human accountability \cite{src-siam-ai}.
\item The length (19 pages in \texttt{amsart}) is ordinary for SIADS.
\end{itemize}
Risk: an editor may find a purely theoretical note on linear SDEs narrow for SIADS. Section 5 (invariant ellipsoids,
admissible times) and Figure 1 supply the dynamical-systems content; the cover letter should stress them.

\emph{Alternative 1: J.\ Math.\ Anal.\ Appl.\ (JMAA).}
\begin{itemize}
\item The scope covers functional analysis and operator theory, dynamical systems and probability \cite{src-jmaa}.
\item JMAA published the closely related composition-operator result \cite{IIS}.
\item It is a safe venue for a complete, correct result of this size.
\item Elsevier's AI declaration is required \cite{src-elsevier}.
\end{itemize}

\emph{Alternative 2: Proceedings of the AMS.} Papers may not exceed 15 printed pages \cite{src-pams}. The core
(Sections 1--5 and Example 6.1) fits if Appendix A, Remarks 7.3 and 7.4 and Proposition 7.5 with Example 7.6 move to
supplementary material. PAMS gives high visibility to short complete answers, but an editor may consider the topic
specialized.

\emph{Not recommended: PAMM}, where \PS\ appeared. It publishes contributions to the GAMM annual meetings
\cite{src-pswpn}, not independent research articles of this kind.

%======================================================================
\begin{thebibliography}{99}
%======================================================================
\raggedright

\bibitem{Aronszajn}
N.~Aronszajn, \emph{Theory of reproducing kernels}, Trans.\ Amer.\ Math.\ Soc.\ \textbf{68} (1950), 337--404.

\bibitem{CarswellMacCluerSchuster}
B.~J.~Carswell, B.~D.~MacCluer, and A.~Schuster, \emph{Composition operators on the Fock space}, Acta Sci.\ Math.\
(Szeged) \textbf{69} (2003), 871--887.

\bibitem{Hertel}
M.~Hertel, F.~M.~Philipp, M.~Schaller, and K.~Worthmann, \emph{Koopman for stochastic dynamics: error bounds for kernel
extended dynamic mode decomposition (kernel EDMD)}, SIAM J.\ Appl.\ Dyn.\ Syst.\ \textbf{25} (2026), no.~3, 2044--2083;
arXiv:2512.20247.

\bibitem{IIS}
M.~Ikeda, I.~Ishikawa, and Y.~Sawano, \emph{Boundedness of composition operators on reproducing kernel Hilbert spaces
with analytic positive definite functions}, J.\ Math.\ Anal.\ Appl.\ \textbf{511} (2022), no.~1, Paper No.~126048.

\bibitem{Klus}
S.~Klus, I.~Schuster, and K.~Muandet, \emph{Eigendecompositions of transfer operators in reproducing kernel Hilbert
spaces}, J.\ Nonlinear Sci.\ \textbf{30} (2020), 283--315.

\bibitem{Koehne}
F.~K\"ohne, F.~M.~Philipp, M.~Schaller, A.~Schiela, and K.~Worthmann, \emph{$L^\infty$-error bounds for approximations
of the Koopman operator by kernel extended dynamic mode decomposition}, SIAM J.\ Appl.\ Dyn.\ Syst.\ \textbf{24} (2025),
501--529.

\bibitem{Koopman}
B.~O.~Koopman, \emph{Hamiltonian systems and transformation in Hilbert space}, Proc.\ Natl.\ Acad.\ Sci.\ USA
\textbf{17} (1931), 315--318.

\bibitem{Kostic}
V.~Kostic, P.~Novelli, A.~Maurer, C.~Ciliberto, L.~Rosasco, and M.~Pontil, \emph{Learning dynamical systems via Koopman
operator regression in reproducing kernel Hilbert spaces}, Advances in Neural Information Processing Systems
\textbf{35} (2022), 4017--4031.

\bibitem{Mauroy}
A.~Mauroy, I.~Mezi\'c, and Y.~Susuki (eds.), \emph{The Koopman Operator in Systems and Control: Concepts,
Methodologies, and Applications}, Lecture Notes in Control and Information Sciences, vol.~484, Springer, Cham, 2020.

\bibitem{PaulsenRaghupathi}
V.~I.~Paulsen and M.~Raghupathi, \emph{An Introduction to the Theory of Reproducing Kernel Hilbert Spaces}, Cambridge
Studies in Advanced Mathematics, vol.~152, Cambridge University Press, Cambridge, 2016.

\bibitem{PhilippACHA}
F.~Philipp, M.~Schaller, K.~Worthmann, S.~Peitz, and F.~N\"uske, \emph{Error bounds for kernel-based approximations of
the Koopman operator}, Appl.\ Comput.\ Harmon.\ Anal.\ \textbf{71} (2024), 101657.

\bibitem{PSWPN}
F.~Philipp, M.~Schaller, K.~Worthmann, S.~Peitz, and F.~N\"uske, \emph{Invariance of Gaussian RKHSs under Koopman
operators of stochastic differential equations with constant matrix coefficients}, Proc.\ Appl.\ Math.\ Mech.\
(PAMM) \textbf{25} (2025), no.~1, e202400127; arXiv:2405.14429 (cited here as \PS).

\bibitem{SHS}
I.~Steinwart, D.~Hush, and C.~Scovel, \emph{An explicit description of the reproducing kernel Hilbert spaces of
Gaussian RBF kernels}, IEEE Trans.\ Inform.\ Theory \textbf{52} (2006), 4635--4643.

\bibitem{VanLoan}
C.~F.~Van~Loan, \emph{Computing integrals involving the matrix exponential}, IEEE Trans.\ Automat.\ Control
\textbf{23} (1978), 395--404.

\bibitem{Vatiwutipong}
P.~Vatiwutipong and N.~Phewchean, \emph{Alternative way to derive the distribution of the multivariate
Ornstein--Uhlenbeck process}, Adv.\ Difference Equ.\ (2019), Paper No.~276.

\bibitem{Wendland}
H.~Wendland, \emph{Scattered Data Approximation}, Cambridge Monographs on Applied and Computational Mathematics,
vol.~17, Cambridge University Press, Cambridge, 2005.

\bibitem{WilliamsEDMD}
M.~O.~Williams, I.~G.~Kevrekidis, and C.~W.~Rowley, \emph{A data-driven approximation of the Koopman operator:
extending dynamic mode decomposition}, J.\ Nonlinear Sci.\ \textbf{25} (2015), 1307--1346.

\bibitem{WilliamsKernel}
M.~O.~Williams, C.~W.~Rowley, and I.~G.~Kevrekidis, \emph{A kernel-based method for data-driven Koopman spectral
analysis}, J.\ Comput.\ Dyn.\ \textbf{2} (2015), 247--265.

\bibitem{src-survey}
Koopman operator theory: fundamentals, control, and applications, arXiv:2607.01819 (July 2026). Search result only;
\url{https://arxiv.org/abs/2607.01819}.

\bibitem{src-hertel}
Search results for arXiv:2512.20247 and DOI 10.1137/26M1839689: \url{https://arxiv.org/abs/2512.20247},
\url{https://doi.org/10.1137/26M1839689}.

\bibitem{src-koehne}
Search results for DOI 10.1137/24M1650120: \url{https://epubs.siam.org/doi/10.1137/24M1650120}.

\bibitem{src-iis}
\url{https://www.sciencedirect.com/science/article/pii/S0022247X22000622} and \url{https://arxiv.org/abs/1911.11992}.

\bibitem{src-mauroy}
\url{https://link.springer.com/book/10.1007/978-3-030-35713-9}.

\bibitem{src-pswpn}
\url{https://arxiv.org/abs/2405.14429} and \url{https://onlinelibrary.wiley.com/doi/10.1002/pamm.202400127}.

\bibitem{src-siads}
SIAM Journal on Applied Dynamical Systems, about the journal and scope:
\url{https://epubs.siam.org/journal/siads/about}.

\bibitem{src-siam-ai}
SIAM Publications, Editorial Policy on Artificial Intelligence: \url{https://epubs.siam.org/artificial-intelligence}.

\bibitem{src-siam-news}
SIAM News, \emph{A contract of trust: artificial intelligence usage for SIAM journal submissions}:
\url{https://www.siam.org/publications/siam-news/articles/a-contract-of-trust-artificial-intelligence-usage-for-siam-journal-submissions/}.

\bibitem{src-elsevier}
Elsevier, Generative AI policies for journals:
\url{https://www.elsevier.com/about/policies-and-standards/generative-ai-policies-for-journals}.

\bibitem{src-jmaa}
Journal of Mathematical Analysis and Applications:
\url{https://www.sciencedirect.com/journal/journal-of-mathematical-analysis-and-applications}.

\bibitem{src-pams}
Proceedings of the American Mathematical Society: \url{https://www.ams.org/proc}.

\end{thebibliography}

\end{document}
